% Options for packages loaded elsewhere
\PassOptionsToPackage{unicode}{hyperref}
\PassOptionsToPackage{hyphens}{url}
\documentclass[
]{article}
\usepackage{xcolor}
\usepackage{amsmath,amssymb}
\setcounter{secnumdepth}{-\maxdimen} % remove section numbering
\usepackage{iftex}
\ifPDFTeX
  \usepackage[T1]{fontenc}
  \usepackage[utf8]{inputenc}
  \usepackage{textcomp} % provide euro and other symbols
\else % if luatex or xetex
  \usepackage{unicode-math} % this also loads fontspec
  \defaultfontfeatures{Scale=MatchLowercase}
  \defaultfontfeatures[\rmfamily]{Ligatures=TeX,Scale=1}
\fi
\usepackage{lmodern}
\ifPDFTeX\else
  % xetex/luatex font selection
\fi
% Use upquote if available, for straight quotes in verbatim environments
\IfFileExists{upquote.sty}{\usepackage{upquote}}{}
\IfFileExists{microtype.sty}{% use microtype if available
  \usepackage[]{microtype}
  \UseMicrotypeSet[protrusion]{basicmath} % disable protrusion for tt fonts
}{}
\makeatletter
\@ifundefined{KOMAClassName}{% if non-KOMA class
  \IfFileExists{parskip.sty}{%
    \usepackage{parskip}
  }{% else
    \setlength{\parindent}{0pt}
    \setlength{\parskip}{6pt plus 2pt minus 1pt}}
}{% if KOMA class
  \KOMAoptions{parskip=half}}
\makeatother
\usepackage{color}
\usepackage{fancyvrb}
\newcommand{\VerbBar}{|}
\newcommand{\VERB}{\Verb[commandchars=\\\{\}]}
\DefineVerbatimEnvironment{Highlighting}{Verbatim}{commandchars=\\\{\}}
% Add ',fontsize=\small' for more characters per line
\newenvironment{Shaded}{}{}
\newcommand{\AlertTok}[1]{\textcolor[rgb]{1.00,0.00,0.00}{\textbf{#1}}}
\newcommand{\AnnotationTok}[1]{\textcolor[rgb]{0.38,0.63,0.69}{\textbf{\textit{#1}}}}
\newcommand{\AttributeTok}[1]{\textcolor[rgb]{0.49,0.56,0.16}{#1}}
\newcommand{\BaseNTok}[1]{\textcolor[rgb]{0.25,0.63,0.44}{#1}}
\newcommand{\BuiltInTok}[1]{\textcolor[rgb]{0.00,0.50,0.00}{#1}}
\newcommand{\CharTok}[1]{\textcolor[rgb]{0.25,0.44,0.63}{#1}}
\newcommand{\CommentTok}[1]{\textcolor[rgb]{0.38,0.63,0.69}{\textit{#1}}}
\newcommand{\CommentVarTok}[1]{\textcolor[rgb]{0.38,0.63,0.69}{\textbf{\textit{#1}}}}
\newcommand{\ConstantTok}[1]{\textcolor[rgb]{0.53,0.00,0.00}{#1}}
\newcommand{\ControlFlowTok}[1]{\textcolor[rgb]{0.00,0.44,0.13}{\textbf{#1}}}
\newcommand{\DataTypeTok}[1]{\textcolor[rgb]{0.56,0.13,0.00}{#1}}
\newcommand{\DecValTok}[1]{\textcolor[rgb]{0.25,0.63,0.44}{#1}}
\newcommand{\DocumentationTok}[1]{\textcolor[rgb]{0.73,0.13,0.13}{\textit{#1}}}
\newcommand{\ErrorTok}[1]{\textcolor[rgb]{1.00,0.00,0.00}{\textbf{#1}}}
\newcommand{\ExtensionTok}[1]{#1}
\newcommand{\FloatTok}[1]{\textcolor[rgb]{0.25,0.63,0.44}{#1}}
\newcommand{\FunctionTok}[1]{\textcolor[rgb]{0.02,0.16,0.49}{#1}}
\newcommand{\ImportTok}[1]{\textcolor[rgb]{0.00,0.50,0.00}{\textbf{#1}}}
\newcommand{\InformationTok}[1]{\textcolor[rgb]{0.38,0.63,0.69}{\textbf{\textit{#1}}}}
\newcommand{\KeywordTok}[1]{\textcolor[rgb]{0.00,0.44,0.13}{\textbf{#1}}}
\newcommand{\NormalTok}[1]{#1}
\newcommand{\OperatorTok}[1]{\textcolor[rgb]{0.40,0.40,0.40}{#1}}
\newcommand{\OtherTok}[1]{\textcolor[rgb]{0.00,0.44,0.13}{#1}}
\newcommand{\PreprocessorTok}[1]{\textcolor[rgb]{0.74,0.48,0.00}{#1}}
\newcommand{\RegionMarkerTok}[1]{#1}
\newcommand{\SpecialCharTok}[1]{\textcolor[rgb]{0.25,0.44,0.63}{#1}}
\newcommand{\SpecialStringTok}[1]{\textcolor[rgb]{0.73,0.40,0.53}{#1}}
\newcommand{\StringTok}[1]{\textcolor[rgb]{0.25,0.44,0.63}{#1}}
\newcommand{\VariableTok}[1]{\textcolor[rgb]{0.10,0.09,0.49}{#1}}
\newcommand{\VerbatimStringTok}[1]{\textcolor[rgb]{0.25,0.44,0.63}{#1}}
\newcommand{\WarningTok}[1]{\textcolor[rgb]{0.38,0.63,0.69}{\textbf{\textit{#1}}}}
\setlength{\emergencystretch}{3em} % prevent overfull lines
\providecommand{\tightlist}{%
  \setlength{\itemsep}{0pt}\setlength{\parskip}{0pt}}
\usepackage{bookmark}
\IfFileExists{xurl.sty}{\usepackage{xurl}}{} % add URL line breaks if available
\urlstyle{same}
\hypersetup{
  hidelinks,
  pdfcreator={LaTeX via pandoc}}

\author{}
\date{}

\begin{document}

\section{LeanDojo: Theorem Proving with Retrieval-Augmented Language
Models}\label{leandojo-theorem-proving-with-retrieval-augmented-language-models}

Kaiyu Yang \(^{1}\) , Aidan M. Swope \(^{2}\) , Alex Gu \(^{3}\) , Rahul
Chalamala \(^{1}\) , Peiyang Song \(^{4}\) , Shixing Yu \(^{5}\) , Saad
Godil \(^{*}\) , Ryan Prenger \(^{2}\) , Anima Anandkumar \(^{1,2}\)

\(^{1}\) Caltech, \(^{2}\) NVIDIA, \(^{3}\) MIT, \(^{4}\) UC Santa
Barbara, \(^{5}\) UT Austin https://leandojo.org

\section{Abstract}\label{abstract}

Large language models (LLMs) have shown promise in proving formal
theorems using proof assistants such as Lean. However, existing methods
are difficult to reproduce or build on, due to private code, data, and
large compute requirements. This has created substantial barriers to
research on machine learning methods for theorem proving. This paper
removes these barriers by introducing LeanDojo: an open-source Lean
playground consisting of toolkits, data, models, and benchmarks.
LeanDojo extracts data from Lean and enables interaction with the proof
environment programmatically. It contains fine-grained annotations of
premises in proofs, providing valuable data for premise selection---a
key bottleneck in theorem proving. Using this data, we develop ReProver
(Retrieval-Augmented Prover): an LLM-based prover augmented with
retrieval for selecting premises from a vast math library. It is
inexpensive and needs only one GPU week of training. Our retriever
leverages LeanDojo's program analysis capability to identify accessible
premises and hard negative examples, which makes retrieval much more
effective. Furthermore, we construct a new benchmark consisting of
98,734 theorems and proofs extracted from Lean's math library. It
features challenging data split requiring the prover to generalize to
theorems relying on novel premises that are never used in training. We
use this benchmark for training and evaluation, and experimental results
demonstrate the effectiveness of ReProver over non-retrieval baselines
and GPT-4. We thus provide the first set of open-source LLM-based
theorem provers without any proprietary datasets and release it under a
permissive MIT license to facilitate further research.

\section{1 Introduction}\label{introduction}

Reasoning is a cornerstone of human intelligence and a fundamental goal
of AI {[}3{]}. One prominent task is automated theorem proving (ATP):
automatically generating proofs for theorems expressed in formal logic.
ATP is useful for formal mathematics, producing mathematical proofs that
can be checked rigorously {[}4{]}. Furthermore, it underpins formal
verification, which is essential for proving the correctness and safety
of high-stakes applications {[}5, 6{]}.

ATP is challenging since the search space is prohibitively large. In
many applications, it is impractical to generate proofs fully
automatically. Therefore, interactive theorem proving (ITP) has emerged
as an alternative paradigm. In ITP, proofs are constructed by human
experts interacting with software tools called proof assistants, such as
Coq {[}7{]}, Isabelle {[}8{]}, and Lean {[}1{]}. Machine learning can
automate such interactive theorem proving, opening up a new avenue for
theorem proving {[}9{]}. The model can learn to interact with proof
assistants, given data containing human-written proofs.

\textit{[Image omitted: images/5e346e38cd6f88e635d6bcd3bf3e5020ac78b1d997f17acb1326d9b8a525ab93.jpg]}

flowchart

Machine learning workflow diagram showing proof tree, local context,
machine learning model, and accessible premises for machine learning.

Figure 1: Top right: LeanDojo extracts proofs in Lean {[}1{]} into
datasets for training machine learning models. It also enables the
trained model to prove theorems by interacting with Lean's proof
environment. Top left: The proof tree of a Lean theorem
\(\forall n \in \mathbb{N}\) , gcd n n = n, where gcd is the greatest
common divisor (details in Sec. 3). When proving the theorem, we start
from the original theorem as the initial state (the root) and repeatedly
apply tactics (the edges) to decompose states into simpler sub-states,
until all states are solved (the leaf nodes). Tactics may rely on
premises such as mod\_self and gcd\_zero\_left defined in a large math
library. E.g., mod\_self is an existing theorem
\(\forall n \in \mathbb{N}\) , n \% n = 0 used in the proof to simplify
the goal. Bottom: Our ReProver model (Sec. 5). Given a state, it
retrieves premises from the math library, which are concatenated with
the state and fed into an encoder-decoder Transformer {[}2{]} to
generate the next tactic.

Formal theorem proving serves as an important challenge for machine
learning. From a computer science perspective, formal proofs can be
treated as programs {[}10{]}. But unlike conventional programs in C++ or
Python, the correctness of proofs can be verified using proof
assistants. Therefore, theorem proving may be considered a special form
of code generation, with rigorous evaluation and no room for the model
to hallucinate. This can be consequential to current large language
models (LLMs), as they have demonstrated exceptional capability in code
generation {[}11{]} but have flaws in factuality and hallucination
{[}12{]}. In addition, augmenting LLMs with external tools, such as
proof assistants, has shown promise in improving their various
capabilities, including multi-step reasoning {[}13{]}.

Current research on LLMs for theorem proving is facing many barriers. To
our knowledge, none of the existing LLM-based provers are open-source
{[}14--21{]}. They all use private pretraining data, and the compute
requirements can reach thousands of GPU days {[}17{]}. Furthermore, some
rely on tailored infrastructure for distributed training and interaction
with the proof assistant---both are not possible to fully reproduce
without open-source code {[}17, 19{]}. We change the status quo by
introducing LeanDojo: open-source toolkits, models, and benchmarks that
give researchers access to state-of-the-art LLM-based provers with
modest computational costs.

Tools for Data Extraction and Interaction. We focus on Lean, a proof
assistant popular among mathematicians. \(^{2}\) Our framework LeanDojo
provides two essential functions for learning-based theorem proving
(Fig. 1): extracting data and enabling models to interact with Lean
programmatically.

For data extraction, LeanDojo extracts training data not directly
visible in the raw Lean code (Fig. 2), e.g., proof trees consisting of
intermediate states between proof steps (Fig. 1 Top left). In addition,
LeanDojo is the first tool to locate premises in Lean proofs, enabling
training machine learning models for premise selection. For interaction,
LeanDojo turns Lean into a gym-like interactive environment {[}22{]}.
Using LeanDojo, the model can observe proof states, change the state by
executing

proof steps (referred to as ``tactics'' in proof assistants), and
receive feedback from Lean. LeanDojo is the first tool capable of
interacting with Lean reliably, reducing proof-checking errors in
existing tools {[}19{]} (correct proofs misjudged as incorrect) from
21.1\% to 1.4\%.

Retrieval-Augmented LLMs for Theorem Proving. LeanDojo addresses a key
bottleneck in theorem proving: premise selection {[}23, 24{]}. Existing
LLM-based provers generate the next proof step (tactic), taking only the
current state as input. However, proving theorems depends critically on
the premises, such as lemmas and definitions, from a math library.

For example, Fig. 1 (Top left) illustrates the proof of ``
\(\forall n \in N\) , gcd n n = n'', where gcd stands for greatest
common divisor. The proof starts from the original theorem as the
initial state and repeatedly applies tactics to decompose states into
simpler sub-states, until all states are solved. Tactics may rely on
premises such as mod\_self and gcd\_zero\_left defined in a large math
library. E.g., mod\_self is an existing theorem `` \(\forall n \in N\) ,
n \% n = 0'' useful for simplifying the goal.

Incorporating all possible premises is too large to fit into LLMs'
input, given the limited context window. Existing methods must learn to
memorize the association between the proof state and the name mod\_self.
It works if the premise has been used in the training data to solve
similar goals, but does not generalize to truly novel scenarios, e.g.,
theorems requiring lemmas unseen in training.

One potential solution is to complement memorization with explicit
premise selection. LeanDojo extracts premise data from Lean, including
where they are defined and used. It enables us to tackle premise
selection by augmenting LLMs with retrieval. We introduce ReProver
(Retrieval-Augmented Prover) (Fig. 1 Bottom): Given the current state,
it generates a tactic conditioning on a small number of premises
retrieved from Lean's math library, mathlib {[}25{]}.

We need to limit retrieval to a small number of premises for it to be
effective, and ideally, they should contain the ground truth premise.
Our retriever builds upon Dense Passage Retriever (DPR) {[}26{]} but
incorporates two algorithmic innovations: First, not all premises are
accessible when proving a theorem (Sec. 3). LeanDojo can perform program
analysis on Lean code to determine accessible premises. On our data,
that reduces the average number of premises from 128K to 33K,
significantly simplifying the retriever's task. Second, DPR needs
negative examples in training and benefits from hard negatives, i.e.,
irrelevant premises that are hard to distinguish from ground truth ones.
We propose in-file negatives: a simple mechanism to find hard negatives
in premise selection, which samples negative premises defined in the
same Lean source file as the ground truth premise.

LeanDojo Benchmark. Using LeanDojo, we construct a benchmark containing
98,734 theorems/proofs extracted from mathlib. Our benchmark is one of
the largest math-focused theorem-proving datasets. We find that the
common practice of splitting theorems randomly into training/testing has
led to an overestimated performance in the previous papers. LLMs can
prove seemingly difficult theorems simply by memorizing the proofs of
similar theorems during training. In LeanDojo Benchmark, we mitigate
this issue by designing challenging data split requiring the model to
generalize to theorems relying on novel premises that are never used in
training.

We use LeanDojo Benchmark to train and evaluate ReProver. Training takes
only five days on a single GPU. In evaluation, ReProver can prove 51.2\%
theorems, outperforming a baseline that generates tactics directly
without retrieval (47.6\%) and another baseline using GPT-4 {[}27{]} to
generate tactics in a zero-shot manner (29.0\%). We also test ReProver
on two existing datasets, MiniF2F {[}28{]} and ProofNet {[}29{]}. It can
prove 26.5\% theorems in MiniF2F and 13.8\% in ProofNet, which is
competitive with state-of-the-art methods without reinforcement learning
{[}19{]}, even though trained using far fewer resources. Moreover, it
can prove 65 theorems that currently do not have proofs in Lean. Thus,
our tool can also serve as an effective tool for augmenting existing
math libraries in Lean.

Contributions. In summary, we make four main contributions: First, we
introduce tools for extracting data from and interacting with Lean.
Second, we develop ReProver, the first retrieval-augmented language
model for theorem proving. Third, we construct a challenging benchmark
for learning-based theorem proving and use it to validate the
effectiveness of ReProver. Finally, we facilitate open research on LLMs
for theorem proving by releasing our data, model, and code. Our method
does not rely on private datasets and can be trained on a single GPU
within a week. We believe this will significantly lower the barriers to
academic research in this area and establish the first accessible
baselines for future work to build upon. Further, our method can be used
to automatically generate new Lean proofs without requiring human
effort.

\section{2 Related Work}\label{related-work}

Theorem Proving. Classical provers express theorems in first-order logic
and search for proofs automatically in a large space {[}30, 31{]}. Even
with data-driven search heuristics {[}32, 33{]}, they fail to scale to
large formalization projects. Therefore, recent work on learning-based
theorem proving has focused on an alternative paradigm: automating the
interaction with proof assistants.

The architecture of learning-based provers progressed from classical
machine learning algorithms such as KNN {[}34{]}, to graph neural
networks explicitly encoding the syntax of formal expressions {[}9,
35{]}, and now Transformer-based LLMs treating expressions as plain
strings {[}14{]}. Besides the model architecture, researchers have
explored several complementary dimensions: proof search algorithms for
assembling model-generated steps into complete proofs {[}17, 21{]};
overcoming data scarcity through reinforcement learning (RL) {[}17, 19,
36, 37{]} or synthetic/auxiliary data {[}16, 38--40{]}; as well as
outsourcing some proof goals to classical provers {[}18, 41--43{]}. Our
base model without retrieval is a combination of straightforward design
choices. It generates tactics by finetuning an encoder-decoder
Transformer, ByT5 {[}44{]}, via supervised learning without RL or
auxiliary data. Then it searches for proofs using best-first search. Our
model's algorithmic novelty lies in the retrieval.

Premise Selection. Selecting useful premises is recognized as a key
challenge in theorem proving {[}23, 24, 45, 46{]}. Machine learning
methods for premise selection have also progressed from classical models
{[}41, 47, 48{]}, recurrent neural networks {[}24{]}, graph neural
networks {[}38{]}, to Transformers {[}49, 50{]}. However, existing
methods either tackle premise selection in isolation without theorem
proving {[}24, 38, 48{]} or feed the premises to a symbolic prover
{[}41, 47, 49{]}. To our knowledge, we are the first to augment a
learning-based formal theorem prover with retrieved premises so that the
prover can learn how to use them effectively. For example, it can decide
whether to use an explicitly retrieved premise or an implicitly
memorized one.

Data and Tools for Theorem Proving. Tools for data extraction and
interacting with proof assistants have been crucial drivers of
learning-based theorem proving. Existing tools and datasets can be
divided by proof assistants: Coq has GamePad {[}51{]}, CoqGym {[}9{]},
and PRISM {[}52{]}; Isabelle has IsarStep {[}53{]} and PISA {[}15{]};
HOL Light has HOList {[}54{]} and HoLStep {[}55{]}, and Lean has
LeanStep {[}16{]} and lean-gym {[}19{]}. MiniF2F {[}28{]} is the only
cross-system dataset, with 488 theorems for evaluation. However, it does
not have training theorems and is restricted to the domain of math
olympiads.

Among available tools extracting data from proof assistants, LeanDojo is
the only one that can extract premises for retrieval-augmented theorem
proving. A few existing datasets also have premises {[}49, 54{]}, but
their data extraction tools are not public, making it difficult to
construct new datasets. In addition, LeanDojo is the only tool that can
interact with Lean robustly (Sec. 4) and can extract data from Lean 4.
See Appendix A.3 for a detailed comparison between LeanDojo and
alternatives.

Mathematical Reasoning in Natural Language. We focus on proving theorems
expressed in formal logic, whereas researchers have also produced a
plethora of work on mathematical reasoning in natural language
{[}56--63{]}. A particularly relevant task is autoformalization,
translating natural language texts into formal theorems and proofs
{[}29, 64--72{]}.

Retrieval-Augmented Language Models. Our ReProver is the first
retrieval-augmented language model for formal theorem proving, though
similar architectures have been studied extensively in NLP {[}73--81{]}.
In addition, there have been many retrieval-augmented methods for code
generation {[}82--88{]}. Most of them retrieve from a corpus not
directly related to the current file, e.g., GitHub or Stack Overflow. In
contrast, our retrieval corpus consists of premises accessible to the
current file, which is determined by program analysis using LeanDojo.
This is similar to what CoCoMIC {[}88{]} does for Python. However, their
retrieval is based on heuristics, whereas ours is learned.

\section{3 Background: Theorem Proving in
Lean}\label{background-theorem-proving-in-lean}

At a high level, Lean is a programming language that allows you to write
not only conventional programs but also theorems and proofs. To that
end, it provides two pieces of machinery: First, it provides a unified
language for defining programs, mathematical objects, theorems, and
proofs, based on functional programming with dependent types {[}89{]}.
Second, it provides a tactic system for constructing machine-checkable
proofs semi-automatically.

\textit{[Image omitted: images/e575e56c92edede68c436b4fcf1378c96292ee85faf662a27dc4298bf01f347a.jpg]}

text\_image

data/nat/lemmas.lean theorem mod\_self (n : nat) : n \% n = 0 := begin
rw {[}mod\_eq\_sub\_mod (le\_ref1 \_), nat.sub\_self, zero\_mod{]} end

def gcd : nat → nat → nat -- gcd z y \textbar{} 0 y := y -- Case 1: z ==
0 \textbar{} (x + 1) y := gcd (y \% (x + 1)) (x + 1) -- Case 2: z
\textgreater{} 0

theorem gcd\_zero\_left (x : nat) : gcd 0 x = x := begin simp {[}gcd{]}
end

theorem gcd\_self (n : nat) : gcd n n = n := begin cases n, \{ unfold
gcd \}, unfold gcd, rewrite mod\_self, apply gcd\_zero\_left end

Figure 2: Definition of greatest common divisor (gcd) in Lean and two
related theorems. The proof of gcd\_self (between ``begin'' and ``end'')
relies on a premise mod\_self imported from another file in the math
library. Lean can run this proof to produce the proof tree in Fig.1 (Top
left).

We use a simple example in Fig. 2 to illustrate how theorems are
formalized and proved in Lean. \(^{3}\) Here we want to formalize the
greatest common divisor (gcd) of two natural numbers. First, we define
gcd as a recursive function, taking two natural numbers as parameters
and returning their gcd via the Euclidean algorithm. Then, we state a
lemma named gcd\_zero\_left that \(\forall x \in \mathbb{N}\) , gcd 0 x
= x, which can be proved simply by the definition of gcd. Finally, we
state our main theorem gcd\_self that \(\forall n \in \mathbb{N}\) , gcd
n n = n, followed by its proof consisting of five tactics. In theorem
proving, we are only concerned with generating the proof, i.e., the part
between ``begin'' and ``end''; everything before ``begin'' is known,
including other files imported.

The syntax of tactics is quite expressive. They can take arguments and
can be combined into compound tactics. You can think of tactics as
programs in a domain-specific language (DSL). Users can extend the DSL
by defining new tactics. This discrete, combinatorial, and unbounded
action space makes theorem proving challenging for machine learning.

Another challenge is premise selection. Premises are existing lemmas or
definitions useful for proving a theorem. They are used as arguments in
tactics. For example, in Fig. 2 and Fig. 1 (Top left), the tactic
``rewrite mod\_self'' rewrites the goal using the premise mod\_self,
which is defined in another file imported by the current file. Proofs
cannot use premises that haven't been defined. For example, gcd\_self
cannot be used to prove gcd\_zero\_left. In addition, they cannot use
premises not imported to the current file. Still, premises come from a
large math library containing hundreds of thousands of existing
definitions and theorems, making it hard, for humans and machines alike,
to select the right premises when generating a tactic. This is a key
bottleneck in theorem proving and is what we aim to address through
retrieval-augmented LLMs.

\section{4 LeanDojo: Toolkit and
Benchmark}\label{leandojo-toolkit-and-benchmark}

LeanDojo serves two essential needs of learning-based theorem proving in
Lean. First, it extracts training data from Lean, and we use this
capability to construct a challenging theorem proving benchmark. Second,
it enables the model to interact with Lean programmatically.

Data Extraction. Lean repos (e.g., mathlib or lean-liquid) contain
source code of human-written theorems/proofs. However, the raw code is
unsuitable for training the prover. It lacks runtime information that
humans can access when using Lean, such as intermediate states between
proof steps. Therefore, LeanDojo extracts the following information not
directly visible in the code:

\begin{itemize}
\tightlist
\item
  File dependencies and abstract syntax trees (ASTs): LeanDojo processes
  the repo to produce a directed acyclic graph whose nodes are files and
  edges are import relations between files. In addition, LeanDojo
  produces the AST of each file. File dependencies and ASTs are useful
  for program analysis, e.g., collecting theorems defined in a file or
  premises accessible to a theorem.\\
\item
  States and tactics: LeanDojo extracts all tactics in proofs. For each
  tactic, it also extracts the states before/after the tactic, which
  allows us to reconstruct the proof tree in Fig. 1 (Top left).\\
\item
  Premises: For each premise, such as mod\_self in Fig. 2, LeanDojo
  records where it is defined (location in data/nat/lemma.lean) and
  where it is used (locations across many files). In addition, premises
  have unique fully qualified names (e.g., nat.mod\_self) but are often
  used by ambiguous short names (mod\_self), relying on Lean to perform
  name resolution. LeanDojo is capable of recording their full names.
\end{itemize}

Lean has basic support for exporting dependencies, ASTs, states, and
tactics. However, it cannot resolve the premises' full names and locate
their definitions. Therefore, we modify Lean to record this information
(details in Appendix A.1). The modified Lean is used only for data
extraction but not for evaluation, so we do not risk accidentally
breaking Lean's logical soundness.

LeanDojo Benchmark. We construct a benchmark for premise selection and
theorem proving, named LeanDojo Benchmark. The data is extracted from
mathlib, \(^{4}\) Lean's centralized math library covering diverse
topics such as analysis, algebra, and geometry. \(^{5}\) LeanDojo
Benchmark is one of the largest math-focused theorem proving datasets,
consisting of 98,734 theorems from 3,384 Lean files. Unlike existing
datasets in Lean {[}16{]}, LeanDojo Benchmark also contains the
definitions of 130,262 premises, including not only theorems but also
other definitions that can be used as premises (e.g., gcd in Fig. 2.
Furthermore, the dataset has 217,776 tactics, 129,243 of them with at
least one premise. The average number of premises is 2.13 among tactics
with premises. Appendix B contains additional information on data
format, datasheet {[}90{]}, hosting, and licensing.

\begin{Shaded}
\begin{Highlighting}[]
\NormalTok{src/algebra/quaternion.lean}
\NormalTok{lemma conj\_mul : (a * b).conj = b.conj * a.conj := begin}
\NormalTok{    ext; simp; ring\_exp}
\NormalTok{end}
\NormalTok{lemma conj\_conj\_mul : (a.conj * b).conj = b.conj * a := begin}
\NormalTok{    rw [conj\_mul, conj\_conj]}
\NormalTok{end}
\NormalTok{lemma conj\_mul\_conj : (a * b.conj).conj = b * a.conj := begin}
\NormalTok{    rw [conj\_mul, conj\_conj]}
\NormalTok{end }
\end{Highlighting}
\end{Shaded}

Figure 3: Similar theorems/proofs are common. If splitting them randomly
into training/testing, the model can prove testing theorems by
memorization.

LeanDojo Benchmark has 94,734/2,000/2,000 theorems for
training/validation/testing. It features a challenging data split for
testing the prover's generalization in more realistic scenarios.
Splitting theorems randomly can overestimate the prover's performance,
by allowing it to prove many theorems through memorization. In
human-written Lean code, a common idiom is to have a block of similar
theorems/proofs for slightly different properties of the same math
concept. For example, in Fig. 3, the last two theorems not only look
similar but have identical proofs. If one of them is in training, the
model can easily prove the other one by memorization. This shortcut
enables the model to prove seemingly nontrivial theorems, including
those requiring premises to prove.

To mitigate this issue, besides the random split, we create a
challenging data split named novel\_premises. It requires testing proofs
to use at least one premise that has never been used in training. For
example, the last two theorems in Fig. 3 both use the premise conj\_mul.
If one theorem is in the training set of the novel\_premises split, the
other one must also be in training.

Interacting with Lean. Another important function of LeanDojo is to
interact with Lean programmatically. It turns Lean into a gym-like
environment {[}22{]}, in which the prover can observe the proof state,
run tactics to change the state, and receive feedback on errors or on
proof completion. This environment is indispensable for
evaluating/deploying the prover or training it through RL.

Below is LeanDojo's main interface for interacting with Lean through
tactics. Lean also supports other proof styles not based on tactics.
Although we only support tactic-style proofs, they are sufficiently
general since any proof can be converted to a tactic-style proof.
\(^{6}\)

\begin{itemize}
\tightlist
\item
  initialize(theorem): Given the theorem to prove, LeanDojo returns the
  initial state. A valid state is a string representing current proof
  goals and local contexts (see the nodes in Fig. 1 Top left). When
  there are multiple goals, their strings are concatenated.\\
\item
  run\_tac(state, tactic): Run a tactic on a given state and return the
  next state. The returned state will be an error state if the tactic
  execution is not successful, e.g., due to timeout or inapplicable
  tactic. If the input state is an error, the result can only be an
  error.
\end{itemize}

Building this environment is technically challenging, as Lean is
designed for human users, not machines. LeanDojo is the first tool that
can interact with Lean reliably. Existing tool {[}19{]} is limited:
21.1\% of the ground truth proofs are misjudged as incorrect, due to
issues with how they construct the proof environment, which distorts the
reported performance and produces unreliable feedback when used in
reinforcement learning. In contrast, LeanDojo reduces the number of
misjudgments to 1.4\%. Details are in Appendix A.2.

\section{5 ReProver: Retrieval-Augmented Theorem
Prover}\label{reprover-retrieval-augmented-theorem-prover}

We develop the ReProver model that uses retrieval to select premises
explicitly. At its core is a retrieval-augmented tactic generator (Fig.
1 Bottom). Given the current proof state, it retrieves a handful of
potentially useful premises and generates a tactic conditioning on the
concatenation of the state and retrieved premises. When proving
theorems, the model generates multiple tactic candidates at each step,
which are used in a standard best-first search algorithm to find proofs
{[}16, 18, 19, 28{]}.

Premise Retrieval. Our retriever is based on Dense Passage Retriever
{[}26{]}. Given a state s as the query and a library of candidate
premises \(P = \{p_i\}_{i=1}^N\) , it retrieves a ranked list of m
premises \(\{p_i'\}_{i=1}^m\) from P. In DPR, s and \(p_i\) are both raw
texts but are embedded in a vector space, and we retrieve the top m
premises maximizing the cosine similarity between the state and the
premise.

More formally, we have a function \(f\) parameterized by \(\theta\) for
embedding both the state and the premises into a \(h\) -dimensional
vector space: \(f(s, \theta), f(p_i, \theta) \in \mathbb{R}^h\) . We
retrieve premises maximizing
\(f(s, \theta)^T f(p_i, \theta) / (\|f(s, \theta)\|_2 \|f(p_i, \theta)\|_2)\)
. We choose \(f\) to be a Transformer encoder {[}2{]} followed by
average pooling:
\(f(\cdot, \theta) = \text{AvgPool}(\text{Enc}(\cdot, \theta))\) .

The retrieval is efficient. The premise embeddings \(f(p_i, \theta)\)
can be pre-computed, and we only need one forward pass to compute
\(f(s, \theta)\) . We do not rerank the retrieved premises as in
Magnushammer {[}49{]}, which is more costly since it requires a separate
forward pass for each retrieved premise.

Similar to DPR, we train the retriever by minimizing a contrastive loss
between positive premises and in-batch negative premises. Specifically,
suppose we have a batch of b states. For each state, we sample a
positive premise from the ground truth and n negative premises from
\(P^{7}\) . They are called ``in-batch'' negatives because they are
shared by all states in the batch---Every state is associated with all
\(b \cdot (n + 1)\) premises; at least 1 of them is positive. Let
\(l_{ij} \in \{0, 1\}\) denote whether a state-premise pair
\((s_i, p_j)\) is positive. We minimize the mean squared loss:

\[
\mathcal {L} (\theta) = \sum_ {i = 1} ^ {b} \sum_ {j = 1} ^ {b \cdot (n + 1)} \left| l _ {i j} - \frac {f \left(s _ {i} , \theta\right) ^ {T} f \left(p _ {j} , \theta\right)}{\| f \left(s _ {i} , \theta\right) \| _ {2} \| f \left(p _ {j} , \theta\right) \| _ {2}} \right| ^ {2}. \tag {1}
\]

Retrieving from Accessible Premises. We incorporate into DPR two
insights tailored to premise selection. First, instead of retrieving
from all premises in the math library, we restrict to premises
accessible to the current theorem. They include premises defined in the
same file before the theorem, as well as those imported from other
files. We compute accessible premises for each theorem, relying on
LeanDojo's capability in program analysis (Sec. 4). Focusing on
accessible premises makes P much smaller. LeanDojo Benchmark contains
130,262 premises in total, but the average number of accessible premises
is only 33,160.

In-file Negative Examples. DPR's performance depends critically on the
quality of negative examples {[}91, 92{]}. In early experiments, we
sampled all \(n\) negative premises randomly, and the model often
mistakenly retrieved other premises from the same file as the positive
one. Therefore, we propose a scheme that samples \(k\) in-file negatives
and \(n - k\) random negatives for training.

Tactic Generation. As in Fig. 1 (Bottom), retrieved premises are
concatenated with the state. \(^{8}\) Then an encoder-decoder
Transformer, ByT5 {[}44{]}, takes them as input and generates the
tactic. The model is trained to minimize the cross entropy loss w.r.t.
human-written tactics.

Training ReProver takes substantially less compute than prior methods
(120 GPU hours vs.~more than 1000 hours {[}16, 17{]}). All existing
LLM-based provers pretrain on datasets specific to math and coding
{[}14--20{]}. The pretraining is computationally expensive, and the
datasets are kept private. In contrast, we choose to avoid
domain-specific pretraining and build upon google/byt5-small---a model
checkpoint that is generic, publicly available, and relatively small
(299M parameters vs.~837M {[}16{]} or 600M {[}17{]}). We could see
further benefits from domain-specific pretraining, as in Minerva
{[}57{]}, or stronger LLMs like LLaMA {[}93{]} or StarCoder {[}94{]},
but that is beyond our scope. In addition, our model is finetuned on
human-written tactics only, without auxiliary data {[}16{]} or data
collected through online interaction with Lean {[}17, 19{]}. These
orthogonal directions are valuable but will significantly increase the
method's complexity and compute requirements.

\section{6 Experiments}\label{experiments}

We evaluate ReProver on LeanDojo Benchmark. It outperforms baselines on
premise selection and theorem proving, demonstrating the promise of
theorem proving with retrieval-augmented language models. Experimental
details and hyperparameters are in Appendix C.1.

Premise Selection. For premise selection, we only use tactics in
LeanDojo Benchmark that have at least one premise. The model, based on a
ByT5 encoder, uses the state before a tactic as the query to retrieve
100 premises. Then, we calculate standard metrics in information
retrieval: R@k (recall for the top k retrieved premises) and MRR (mean
reciprocal rank).

Our first baseline is a classical BM25 retriever {[}95{]} without
machine learning. Results in Table 1 show that our method outperforms
BM25 significantly across the board. However, it exhibits a large
performance degradation on the challenging data split (comparing
novel\_premises to random). This is consistent with the general
observation that machine learning can be brittle in the presence of
distribution shifts. In addition, we compare with two ablations: one
retrieving from all premises (instead of accessible premises only) and
the other without in-file negatives. They perform worse than our method,
demonstrating the effectiveness of our two improvements upon DPR.

Theorem Proving Experimental Setup. Then we evaluate ReProver on theorem
proving. The training has two stages: First, we train the retriever and
use it to retrieve 100 premises for all proof states in LeanDojo
Benchmark. Second, we train the tactic generator, taking as input the
concatenation of the state and retrieved premises (truncated to a length
limit). During evaluation, the tactic generator is combined with
best-first search to prove theorems. We evaluate the Pass@1 metric: The
prover is given only one attempt and must find the proof within a wall
time limit of 10 minutes. Training takes five days on a single NVIDIA
A100 GPU with 80GB memory, and evaluation takes two days on eight V100
GPUs. Please see Appendix C.1 for details.

Baselines. Following prior work {[}16, 28{]}, we include tidy as a
baseline. It is a tactic in mathlib that tries to complete the proof
using heuristics (without machine learning). We apply tidy directly

Table 1: Premise selection testing performance. For each method, we
train and evaluate two models independently using different data splits
(random and novel\_premises; see Sec. 4). R@k is the recall for the top
k retrieved premises, and MRR is the mean reciprocal rank metric (higher
is better). Our retriever outperforms BM25 and ablations. Results for
Lean 4 are in Appendix D.

Method

random

novel\_premises

R@1

R@10

MRR

R@1

R@10

MRR

BM25

6.7

17.2

0.15

5.9

15.5

0.14

w/ all premises

1.9

11.9

0.08

2.1

12.4

0.08

Ours

13.5

38.4

0.31

9.1

27.6

0.24

w/ all premises

11.7

36.2

0.27

7.1

23.1

0.20

w/o in-file negatives

10.8

33.1

0.25

7.9

25.7

0.22

to the original theorem and see if it can succeed within the wall time
limit. Another baseline uses GPT-4 as the tactic generator. Given a
state, it queries GPT-4 to generate 35 tactics in zero-shot. After
removing invalid ones, the remaining tactics are combined with
best-first search to find proofs. Data contamination is possible: Many
proofs had been publicly available on GitHub before GPT-4's data cutoff
date (September 2021). See Appendix C.2 for details.

Unfortunately, it is not feasible to compare with existing LLM-based
provers in Lean {[}16, 17, 19{]}. None of them are open-source or can be
reproduced with reasonable effort. Furthermore, we cannot compare
directly with the numbers reported in their papers, due to differences
in data, infrastructure, and training procedures (details in Appendix
C.3). Many difficulties are due to the private nature of existing
methods. By releasing our code and models, we hope to create accessible
baselines for future work to build upon.

Table 2: Theorem proving Pass@1 (\%) on the testing data of LeanDojo
Benchmark. Our ReProver model outperforms tidy, GPT-4, and a baseline
that generates tactics directly without retrieval. Results for Lean 4
are in Appendix D.

Method

random

novel\_premises

tidy

23.8

5.3

GPT-4

29.0

7.4

ReProver (ours)

51.2

26.3

w/o retrieval

47.6

23.2

Results. Table 2 shows the results on the testing data of LeanDojo
Benchmark. ReProver outperforms all baselines on two different data
splits, demonstrating the effectiveness of retrieval-augmented theorem
proving. GPT-4 performs substantially worse than our method, even though
it may have seen the ground truth proofs due to data contamination. The
task cannot be solved out of the box by state-of-the-art LLMs, calling
for algorithmic innovations to make further progress.

Testing theorems in novel\_premises are indeed much more challenging.
All methods in Table 2 perform substantially worse on novel\_premises
than the random split. We argue that performance on challenging splits
is more indicative of the prover's capability and should be emphasized
in the future development of theorem proving.

Evaluation on MiniF2F and ProofNet. We run ReProver to prove theorems in
MiniF2F {[}28{]} and ProofNet {[}29{]}. These two datasets are for
testing only and do not have training theorems, which makes them
challenging since the distribution of theorems is quite different from
mathlib used to train ReProver. MiniF2F focuses on math olympiads, and
ProofNet focuses on exercises in undergraduate math textbooks. On
MiniF2F's test set in Lean, ReProver achieves a Pass@1 of 26.5\%, which
is competitive with state-of-the-art methods without RL (25.9\% in Polu
et al.~{[}19{]}). On ProofNet, our Pass@1 is 13.8\%, which is the first
reported theorem proving result on this dataset. Further, many theorems
do not have ground truth proofs in Lean. Our prover discovers 33 proofs
in MiniF2F and 39 proofs in ProofNet that currently do not have Lean
proofs. Please see Appendix C.4 for details, examples, and caveats.

\section{7 Conclusion}\label{conclusion}

We have introduced LeanDojo: an open-source playground for
learning-based theorem proving in Lean, consisting of toolkits, models,
and benchmarks. It extracts data from Lean and enables the model to
interact with Lean programmatically. We have developed ReProver, the
first retrieval-augmented LLM for theorem proving. Limitations and
future work are discussed in Appendix F.

We have released our code, data, models, and documentation to facilitate
future research:

\begin{itemize}
\tightlist
\item
  LeanDojo's codebase for data extraction and interaction with Lean:
  https://github.com/lean-dojo/LeanDojo\\
\item
  LeanDojo's documentation: https://leandojo.readthedocs.io\\
\item
  Datasets: (1) LeanDojo Benchmark:
  https://doi.org/10.5281/zenodo.8016385 with DOI
  10.5281/zenodo.8016385. (2) LeanDojo Benchmark 4 (Appendix D):
  https://doi.org/10.5281/zenodo.8040109 with DOI
  10.5281/zenodo.8040109.\\
\item
  ReProver's code and models: https://github.com/lean-dojo/ReProver\\
\item
  ChatGPT plugin (Appendix E):
  https://github.com/lean-dojo/LeanDojoChatGPT\\
\item
  LeanDojo Website: https://leandojo.org
\end{itemize}

\section{Acknowledgments and Disclosure of
Funding}\label{acknowledgments-and-disclosure-of-funding}

This work is partially supported by Caltech's Center for Autonomous
Systems and Technologies. Kaiyu Yang is supported by the Computing,
Data, and Society Postdoctoral Fellowship at Caltech. Alex Gu is
supported by the National Science Foundation (NSF) Graduate Research
Fellowship. Rahul Chalamala and Peiyang Song are supported by the Summer
Undergraduate Research Fellowships (SURF) program at Caltech. Anima
Anandkumar is partially supported by the Bren endowed chair. We
appreciate the valuable feedback from Logan Murphy and members of the
Anima AI+Science Lab on an initial version of this paper. We thank
Junyan Xu for manually inspecting the proofs generated by our model on
ProofNet. We also thank Jeremy Avigad and Mario Carneiro for insightful
discussions on supporting Lean 4 in LeanDojo.

\section{References}\label{references}

{[}1{]} Leonardo de Moura, Soonho Kong, Jeremy Avigad, Floris Van Doorn,
and Jakob von Raumer. The Lean theorem prover (system description). In
International Conference on Automated Deduction (CADE), 2015. 1, 2, 22\\
{[}2{]} Ashish Vaswani, Noam Shazeer, Niki Parmar, Jakob Uszkoreit,
Llion Jones, Aidan N Gomez, Łukasz Kaiser, and Illia Polosukhin.
Attention is all you need. In Neural Information Processing Systems
(NeurIPS), 2017. 2, 7\\
{[}3{]} Allen Newell and Herbert Simon. The logic theory machine---a
complex information processing system. IRE Transactions on information
theory, 2(3):61--79, 1956. 1\\
{[}4{]} Kevin Buzzard. The future of mathematics. CRNS-Imperial Lecture,
2019. 1\\
{[}5{]} Xavier Leroy, Sandrine Blazy, Daniel Kästner, Bernhard Schommer,
Markus Pister, and Christian Ferdinand. CompCert---a formally verified
optimizing compiler. In Embedded Real Time Software and Systems, 2016.
1\\
{[}6{]} Talia Ringer, Karl Palmskog, Ilya Sergey, Milos Gligoric,
Zachary Tatlock, et al.~QED at large: A survey of engineering of
formally verified software. Foundations and Trends® in Programming
Languages, 2019. 1\\
{[}7{]} Bruno Barras, Samuel Boutin, Cristina Cornes, Judicaël Courant,
Jean-Christophe Filliatre, Eduardo Gimenez, Hugo Herbelin, Gerard Huet,
Cesar Munoz, Chetan Murthy, et al.~The Coq proof assistant reference
manual: Version 6.1. PhD thesis, Inria, 1997. 1\\
{[}8{]} Tobias Nipkow, Markus Wenzel, and Lawrence C Paulson.
Isabelle/HOL: a proof assistant for higher-order logic. 2002. 1\\
{[}9{]} Kaiyu Yang and Jia Deng. Learning to prove theorems via
interacting with proof assistants. In International Conference on
Machine Learning (ICML), 2019. 1, 4, 35\\
{[}10{]} William A Howard. The formulae-as-types notion of construction.
To HB Curry: Essays on Combinatory Logic, Lambda Calculus and Formalism,
1980. 2\\
{[}11{]} Mark Chen, Jerry Tworek, Heewoo Jun, Qiming Yuan, Henrique
Ponde de Oliveira Pinto, Jared Kaplan, Harri Edwards, Yuri Burda,
Nicholas Joseph, Greg Brockman, et al.~Evaluating large language models
trained on code. arXiv preprint arXiv:2107.03374, 2021. 2\\
{[}12{]} Ziwei Ji, Nayeon Lee, Rita Frieske, Tiezheng Yu, Dan Su, Yan
Xu, Etsuko Ishii, Ye Jin Bang, Andrea Madotto, and Pascale Fung. Survey
of hallucination in natural language generation. ACM Computing Surveys,
55(12):1--38, 2023. 2\\
{[}13{]} Timo Schick, Jane Dwivedi-Yu, Roberto Dessì, Roberta Raileanu,
Maria Lomeli, Luke Zettlemoyer, Nicola Cancedda, and Thomas Scialom.
Toolformer: Language models can teach themselves to use tools. arXiv
preprint arXiv:2302.04761, 2023. 2\\
{[}14{]} Stanislas Polu and Ilya Sutskever. Generative language modeling
for automated theorem proving. arXiv preprint arXiv:2009.03393, 2020. 2,
4, 8\\
{[}15{]} Albert Qiaochu Jiang, Wenda Li, Jesse Michael Han, and Yuhuai
Wu. LISA: Language models of ISAbelle proofs. In Conference on
Artificial Intelligence and Theorem Proving (AITP), 2021. 4\\
{[}16{]} Jesse Michael Han, Jason Rute, Yuhuai Wu, Edward Ayers, and
Stanislas Polu. Proof artifact co-training for theorem proving with
language models. In International Conference on Learning Representations
(ICLR), 2022. 4, 6, 7, 8, 9, 19, 20, 26, 36\\
{[}17{]} Guillaume Lample, Timothee Lacroix, Marie-Anne Lachaux,
Aurelien Rodriguez, Amaury Hayat, Thibaut Lavril, Gabriel Ebner, and
Xavier Martinet. HyperTree proof search for neural theorem proving. In
Neural Information Processing Systems (NeurIPS), 2022. 2, 4, 8, 9, 26,
36\\
{[}18{]} Albert Qiaochu Jiang, Wenda Li, Szymon Tworkowski, Konrad
Czechowski, Tomasz Odrzygóźdź, Piotr Miłoś, Yuhuai Wu, and Mateja
Jamnik. Thor: Wielding hammers to integrate language models and
automated theorem provers. In Neural Information Processing Systems
(NeurIPS), 2022. 4, 7, 26\\
{[}19{]} Stanislas Polu, Jesse Michael Han, Kunhao Zheng, Mantas Baksys,
Igor Babuschkin, and Ilya Sutskever. Formal mathematics statement
curriculum learning. In International Conference on Learning
Representations (ICLR), 2023. 2, 3, 4, 7, 8, 9, 18, 19, 20, 26, 36

{[}20{]} Emily First, Markus N Rabe, Talia Ringer, and Yuriy Brun.
Baldur: Whole-proof generation and repair with large language models.
arXiv preprint arXiv:2303.04910, 2023. 8, 35\\
{[}21{]} Haiming Wang, Ye Yuan, Zhengying Liu, Jianhao Shen, Yichun Yin,
Jing Xiong, Enze Xie, Han Shi, Yujun Li, Lin Li, et al.~DT-Solver:
Automated theorem proving with dynamic-tree sampling guided by
proof-level value function. In Annual Meeting of the Association for
Computational Linguistics (ACL), 2023. 2, 4\\
{[}22{]} Greg Brockman, Vicki Cheung, Ludwig Pettersson, Jonas
Schneider, John Schulman, Jie Tang, and Wojciech Zaremba. OpenAI Gym.
arXiv preprint arXiv:1606.01540, 2016. 2, 7\\
{[}23{]} Josef Urban. MPTP---motivation, implementation, first
experiments. Journal of Automated Reasoning, 33:319--339, 2004. 3, 4\\
{[}24{]} Geoffrey Irving, Christian Szegedy, Alexander A Alemi, Niklas
Eén, François Chollet, and Josef Urban. DeepMath---deep sequence models
for premise selection. In Neural Information Processing Systems
(NeurIPS), 2016. 3, 4\\
{[}25{]} The mathlib Community. The Lean mathematical library. In
Proceedings of the 9th ACM SIGPLAN International Conference on Certified
Programs and Proofs, CPP 2020, pages 367--381, New York, NY, USA, 2020.
Association for Computing Machinery. ISBN 9781450370974. doi:
10.1145/3372885.3373824. URL https://doi.org/10.1145/3372885.3373824.
3\\
{[}26{]} Vladimir Karpukhin, Barlas Oguz, Sewon Min, Patrick Lewis,
Ledell Wu, Sergey Edunov, Danqi Chen, and Wen-tau Yih. Dense passage
retrieval for open-domain question answering. In Conference on Empirical
Methods in Natural Language Processing (EMNLP), 2020. 3, 7, 33\\
{[}27{]} OpenAI. GPT-4 technical report. arXiv preprint
arXiv:2303.08774, 2023. 3, 24, 32\\
{[}28{]} Kunhao Zheng, Jesse Michael Han, and Stanislas Polu. MiniF2F: a
cross-system benchmark for formal olympiad-level mathematics. In
International Conference on Learning Representations (ICLR), 2022. 3, 4,
7, 8, 9, 26, 27\\
{[}29{]} Zhangir Azerbayev, Bartosz Piotrowski, Hailey Schoelkopf,
Edward W Ayers, Dragomir Radev, and Jeremy Avigad. ProofNet:
Autoformalizing and formally proving undergraduate-level mathematics.
arXiv preprint arXiv:2302.12433, 2023. 3, 4, 9, 26, 28, 29\\
{[}30{]} Alan JA Robinson and Andrei Voronkov. Handbook of automated
reasoning, volume 1. 2001. 4\\
{[}31{]} Laura Kovács and Andrei Voronkov. First-order theorem proving
and vampire. In International Conference on Computer Aided Verification
(CAV), 2013. 4\\
{[}32{]} Sarah Loos, Geoffrey Irving, Christian Szegedy, and Cezary
Kaliszyk. Deep network guided proof search. arXiv preprint
arXiv:1701.06972, 2017. 4\\
{[}33{]} James P Bridge, Sean B Holden, and Lawrence C Paulson. Machine
learning for first-order theorem proving: learning to select a good
heuristic. Journal of Automated Reasoning, 53:141--172, 2014. 4\\
{[}34{]} Thibault Gauthier, Cezary Kaliszyk, Josef Urban, Ramana Kumar,
and Michael Norrish. TacticToe: learning to prove with tactics. Journal
of Automated Reasoning, 65:257--286, 2021. 4\\
{[}35{]} Aditya Paliwal, Sarah Loos, Markus Rabe, Kshitij Bansal, and
Christian Szegedy. Graph representations for higher-order logic and
theorem proving. In AAAI Conference on Artificial Intelligence, 2020.
4\\
{[}36{]} Kshitij Bansal, Christian Szegedy, Markus N Rabe, Sarah M Loos,
and Viktor Toman. Learning to reason in large theories without
imitation. arXiv preprint arXiv:1905.10501, 2019. 4\\
{[}37{]} Minchao Wu, Michael Norrish, Christian Walder, and Amir
Dezfouli. TacticZero: Learning to prove theorems from scratch with deep
reinforcement learning. In Neural Information Processing Systems
(NeurIPS), 2021. 4\\
{[}38{]} Mingzhe Wang and Jia Deng. Learning to prove theorems by
learning to generate theorems. In Neural Information Processing Systems
(NeurIPS), 2020. 4, 36\\
{[}39{]} Markus Norman Rabe, Dennis Lee, Kshitij Bansal, and Christian
Szegedy. Mathematical reasoning via self-supervised skip-tree training.
In International Conference on Learning Representations (ICLR), 2021.

{[}40{]} Yuhuai Wu, Markus N Rabe, Wenda Li, Jimmy Ba, Roger B Grosse,
and Christian Szegedy. LIME: Learning inductive bias for primitives of
mathematical reasoning. In International Conference on Machine Learning
(ICML), 2021. 4\\
{[}41{]} Sascha Böhme and Tobias Nipkow. Sledgehammer: judgement day. In
International Joint Conference on Automated Reasoning (IJCAR), 2010. 4\\
{[}42{]} Jasmin Christian Blanchette, Cezary Kaliszyk, Lawrence C
Paulson, and Josef Urban. Hammering towards QED. Journal of Formalized
Reasoning, 9(1):101--148, 2016.\\
{[}43{]} Łukasz Czajka and Cezary Kaliszyk. Hammer for Coq: Automation
for dependent type theory. Journal of Automated Reasoning, 2018. 4\\
{[}44{]} Linting Xue, Aditya Barua, Noah Constant, Rami Al-Rfou, Sharan
Narang, Mihir Kale, Adam Roberts, and Colin Raffel. ByT5: Towards a
token-free future with pre-trained byte-to-byte models. Transactions of
the Association for Computational Linguistics (TACL), 10:291--306, 2022.
4, 8, 24, 32\\
{[}45{]} Christian Szegedy, Markus Rabe, and Henryk Michalewski.
Retrieval-augmented proof step synthesis. In Conference on Artificial
Intelligence and Theorem Proving (AITP), 2021. 4\\
{[}46{]} Yuhuai Wu. Formal premise selection with language models. In
Conference on Artificial Intelligence and Theorem Proving (AITP), 2022.
4\\
{[}47{]} Jesse Alama, Tom Heskes, Daniel Kühlwein, Evgeni Tsivtsivadze,
and Josef Urban. Premise selection for mathematics by corpus analysis
and kernel methods. Journal of Automated Reasoning, 52:191--213, 2014.
4\\
{[}48{]} Bartosz Piotrowski, Ramon Fernández Mir, and Edward Ayers.
Machine-learned premise selection for Lean. In International Conference
on Automated Reasoning with Analytic Tableaux and Related Methods
(TABLEAUX), 2023. 4\\
{[}49{]} Maciej Mikuła, Szymon Antoniak, Szymon Tworkowski, Albert
Qiaochu Jiang, Jin Peng Zhou, Christian Szegedy, Łukasz Kuciński, Piotr
Miłoś, and Yuhuai Wu. Magnushammer: A transformer-based approach to
premise selection. arXiv preprint arXiv:2303.04488, 2023. 4, 7\\
{[}50{]} Eric Yeh, Briland Hitaj, Sam Owre, Maena Quemener, and
Natarajan Shankar. CoProver: A recommender system for proof
construction. arXiv preprint arXiv:2304.10486, 2023. 4\\
{[}51{]} Daniel Huang, Prafulla Dhariwal, Dawn Song, and Ilya Sutskever.
GamePad: A learning environment for theorem proving. In International
Conference on Learning Representations (ICLR), 2019. 4\\
{[}52{]} Tom Reichel, R Henderson, Andrew Touchet, Andrew Gardner, and
Talia Ringer. Proof repair infrastructure for supervised models:
Building a large proof repair dataset. In International Conference on
Interactive Theorem Proving (ITP), 2023. 4\\
{[}53{]} Wenda Li, Lei Yu, Yuhuai Wu, and Lawrence C Paulson. IsarStep:
a benchmark for high-level mathematical reasoning. In International
Conference on Learning Representations (ICLR), 2021. 4\\
{[}54{]} Kshitij Bansal, Sarah Loos, Markus Rabe, Christian Szegedy, and
Stewart Wilcox. HOList: An environment for machine learning of higher
order logic theorem proving. In International Conference on Machine
Learning (ICML), 2019. 4\\
{[}55{]} Cezary Kaliszyk, François Chollet, and Christian Szegedy.
HolStep: A machine learning dataset for higher-order logic theorem
proving. In International Conference on Learning Representations (ICLR),
2017. 4\\
{[}56{]} Dan Hendrycks, Collin Burns, Saurav Kadavath, Akul Arora,
Steven Basart, Eric Tang, Dawn Song, and Jacob Steinhardt. Measuring
mathematical problem solving with the MATH dataset. In Neural
Information Processing Systems (NeurIPS), Datasets and Benchmarks Track,
2021. 4\\
{[}57{]} Aitor Lewkowycz, Anders Johan Andreassen, David Dohan, Ethan
Dyer, Henryk Michalewski, Vinay Venkatesh Ramasesh, Ambrose Slone, Cem
Anil, Imanol Schlag, Theo Gutman-Solo, et al.~Solving quantitative
reasoning problems with language models. In Neural Information
Processing Systems (NeurIPS), 2022. 8\\
{[}58{]} Deborah Ferreira and André Freitas. Premise selection in
natural language mathematical texts. In Annual Meeting of the
Association for Computational Linguistics (ACL), 2020.

{[}59{]} Sean Welleck, Jiacheng Liu, Ronan Le Bras, Hannaneh Hajishirzi,
Yejin Choi, and Kyunghyun Cho. NaturalProofs: Mathematical theorem
proving in natural language. In Neural Information Processing Systems
(NeurIPS), Datasets and Benchmarks Track, 2021.\\
{[}60{]} Sean Welleck, Jiacheng Liu, Ximing Lu, Hannaneh Hajishirzi, and
Yejin Choi. NaturalProver: Grounded mathematical proof generation with
language models. In Neural Information Processing Systems (NeurIPS),
2022.\\
{[}61{]} Swaroop Mishra, Matthew Finlayson, Pan Lu, Leonard Tang, Sean
Welleck, Chitta Baral, Tanmay Rajpurohit, Oyvind Tafjord, Ashish
Sabharwal, Peter Clark, et al.~Lila: A unified benchmark for
mathematical reasoning. In Conference on Empirical Methods in Natural
Language Processing (EMNLP), 2022.\\
{[}62{]} Pan Lu, Liang Qiu, Wenhao Yu, Sean Welleck, and Kai-Wei Chang.
A survey of deep learning for mathematical reasoning. arXiv preprint
arXiv:2212.10535, 2022.\\
{[}63{]} Jordan Meadows and Andre Freitas. A survey in mathematical
language processing. arXiv preprint arXiv:2205.15231, 2022. 4\\
{[}64{]} Qingxiang Wang, Cezary Kaliszyk, and Josef Urban. First
experiments with neural translation of informal to formal mathematics.
In Conferences on Intelligent Computer Mathematics (CICM), 2018. 4\\
{[}65{]} Matthias Cosler, Christopher Hahn, Daniel Mendoza, Frederik
Schmitt, and Caroline Trippel. nl2spec: Interactively translating
unstructured natural language to temporal logics with large language
models. arXiv preprint arXiv:2303.04864, 2023.\\
{[}66{]} Jiayi Pan, Glen Chou, and Dmitry Berenson. Data-efficient
learning of natural language to linear temporal logic translators for
robot task specification. In International Conference on Robotics and
Automation (ICRA), 2023.\\
{[}67{]} Christopher Hahn, Frederik Schmitt, Julia J Tillman, Niklas
Metzger, Julian Siber, and Bernd Finkbeiner. Formal specifications from
natural language. arXiv preprint arXiv:2206.01962, 2022.\\
{[}68{]} Yuhuai Wu, Albert Qiaochu Jiang, Wenda Li, Markus Rabe, Charles
Staats, Mateja Jamnik, and Christian Szegedy. Autoformalization with
large language models. In Neural Information Processing Systems
(NeurIPS), 2022.\\
{[}69{]} Albert Q Jiang, Sean Welleck, Jin Peng Zhou, Wenda Li, Jiacheng
Liu, Mateja Jamnik, Timothée Lacroix, Yuhuai Wu, and Guillaume Lample.
Draft, Sketch, and Prove: Guiding formal theorem provers with informal
proofs. In International Conference on Learning Representations (ICLR),
2023. 26\\
{[}70{]} Xueliang Zhao, Wenda Li, and Lingpeng Kong. Decomposing the
enigma: Subgoal-based demonstration learning for formal theorem proving.
arXiv preprint arXiv:2305.16366, 2023. 26\\
{[}71{]} Garett Cunningham, Razvan C Bunescu, and David Juedes. Towards
autoformalization of mathematics and code correctness: Experiments with
elementary proofs. arXiv preprint arXiv:2301.02195, 2023.\\
{[}72{]} Yongchao Chen, Rujul Gandhi, Yang Zhang, and Chuchu Fan. NL2TL:
Transforming natural languages to temporal logics using large language
models. arXiv preprint arXiv:2305.07766, 2023. 4\\
{[}73{]} Urvashi Khandelwal, Omer Levy, Dan Jurafsky, Luke Zettlemoyer,
and Mike Lewis. Generalization through memorization: Nearest neighbor
language models. In International Conference on Learning Representations
(ICLR), 2020. 4\\
{[}74{]} Kelvin Guu, Kenton Lee, Zora Tung, Panupong Pasupat, and
Mingwei Chang. Retrieval augmented language model pre-training. In
International Conference on Machine Learning (ICML), 2020.\\
{[}75{]} Patrick Lewis, Ethan Perez, Aleksandra Piktus, Fabio Petroni,
Vladimir Karpukhin, Naman Goyal, Heinrich Küttler, Mike Lewis, Wen-tau
Yih, Tim Rocktäschel, et al.~Retrieval-augmented generation for
knowledge-intensive NLP tasks. In Neural Information Processing Systems
(NeurIPS), 2020.\\
{[}76{]} Sebastian Borgeaud, Arthur Mensch, Jordan Hoffmann, Trevor Cai,
Eliza Rutherford, Katie Millican, George Bm Van Den Driessche,
Jean-Baptiste Lespiau, Bogdan Damoc, Aidan Clark, et al.~Improving
language models by retrieving from trillions of tokens. In International
Conference on Machine Learning (ICML), 2022.\\
{[}77{]} Zonglin Li, Ruiqi Guo, and Sanjiv Kumar. Decoupled context
processing for context augmented language modeling. In Neural
Information Processing Systems (NeurIPS), 2022.

{[}78{]} Zhengbao Jiang, Luyu Gao, Jun Araki, Haibo Ding, Zhiruo Wang,
Jamie Callan, and Graham Neubig. Retrieval as attention: End-to-end
learning of retrieval and reading within a single transformer. In
Conference on Empirical Methods in Natural Language Processing (EMNLP),
2022.\\
{[}79{]} Yuhuai Wu, Markus Norman Rabe, DeLesley Hutchins, and Christian
Szegedy. Memorizing transformers. In International Conference on
Learning Representations (ICLR), 2022.\\
{[}80{]} Gautier Izacard, Patrick Lewis, Maria Lomeli, Lucas Hosseini,
Fabio Petroni, Timo Schick, Jane Dwivedi-Yu, Armand Joulin, Sebastian
Riedel, and Edouard Grave. Few-shot learning with retrieval augmented
language models. arXiv preprint arXiv:2208.03299, 2022.\\
{[}81{]} Zexuan Zhong, Tao Lei, and Danqi Chen. Training language models
with memory augmentation. In Conference on Empirical Methods in Natural
Language Processing (EMNLP), 2022. 4\\
{[}82{]} Shirley Anugrah Hayati, Raphael Olivier, Pravalika Avvaru,
Pengcheng Yin, Anthony Tomasic, and Graham Neubig. Retrieval-based
neural code generation. In Conference on Empirical Methods in Natural
Language Processing (EMNLP), 2018. 4\\
{[}83{]} Md Rizwan Parvez, Wasi Ahmad, Saikat Chakraborty, Baishakhi
Ray, and Kai-Wei Chang. Retrieval augmented code generation and
summarization. In Findings of the Association for Computational
Linguistics: EMNLP, 2021.\\
{[}84{]} Shuai Lu, Nan Duan, Hojae Han, Daya Guo, Seung-won Hwang, and
Alexey Svyatkovskiy. ReACC: A retrieval-augmented code completion
framework. In Annual Meeting of the Association for Computational
Linguistics (ACL), 2022.\\
{[}85{]} Shuyan Zhou, Uri Alon, Frank F Xu, Zhengbao Jiang, and Graham
Neubig. DocPrompting: Generating code by retrieving the docs. In
International Conference on Learning Representations (ICLR), 2023.\\
{[}86{]} Disha Shrivastava, Hugo Larochelle, and Daniel Tarlow.
Repository-level prompt generation for large language models of code.
arXiv preprint arXiv:2206.12839, 2022.\\
{[}87{]} Fengji Zhang, Bei Chen, Yue Zhang, Jin Liu, Daoguang Zan, Yi
Mao, Jian-Guang Lou, and Weizhu Chen. RepoCoder: Repository-level code
completion through iterative retrieval and generation. arXiv preprint
arXiv:2303.12570, 2023.\\
{[}88{]} Yangruibo Ding, Zijian Wang, Wasi Uddin Ahmad, Murali Krishna
Ramanathan, Ramesh Nallapati, Parminder Bhatia, Dan Roth, and Bing
Xiang. CoCoMIC: Code completion by jointly modeling in-file and
cross-file context. arXiv preprint arXiv:2212.10007, 2022. 4\\
{[}89{]} David Thrane Christiansen. Functional programming in Lean,
2023. 4\\
{[}90{]} Timnit Gebru, Jamie Morgenstern, Briana Vecchione, Jennifer
Wortman Vaughan, Hanna Wallach, Hal Daumé Iii, and Kate Crawford.
Datasheets for datasets. Communications of the ACM, 64(12):86--92, 2021.
6, 22\\
{[}91{]} Jingtao Zhan, Jiaxin Mao, Yiqun Liu, Jiafeng Guo, Min Zhang,
and Shaoping Ma. Optimizing dense retrieval model training with hard
negatives. In International ACM SIGIR Conference on Research and
Development in Information Retrieval (SIGIR), 2021. 8\\
{[}92{]} Jing Lu, Gustavo Hernandez Abrego, Ji Ma, Jianmo Ni, and Yinfei
Yang. Multi-stage training with improved negative contrast for neural
passage retrieval. In Conference on Empirical Methods in Natural
Language Processing (EMNLP), 2021. 8\\
{[}93{]} Hugo Touvron, Thibaut Lavril, Gautier Izacard, Xavier Martinet,
Marie-Anne Lachaux, Timothée Lacroix, Baptiste Rozière, Naman Goyal,
Eric Hambro, Faisal Azhar, et al.~LLaMA: Open and efficient foundation
language models. arXiv preprint arXiv:2302.13971, 2023. 8\\
{[}94{]} Raymond Li, Loubna Ben Allal, Yangtian Zi, Niklas Muennighoff,
Denis Kocetkov, Chenghao Mou, Marc Marone, Christopher Akiki, Jia Li,
Jenny Chim, et al.~StarCoder: may the source be with you! arXiv preprint
arXiv:2305.06161, 2023. 8, 32\\
{[}95{]} Stephen Robertson, Hugo Zaragoza, et al.~The probabilistic
relevance framework: BM25 and beyond. Foundations and Trends® in
Information Retrieval, 3(4):333--389, 2009. 8\\
{[}96{]} Leonardo de Moura, Jeremy Avigad, Soonho Kong, and Cody Roux.
Elaboration in dependent type theory. arXiv preprint arXiv:1505.04324,
2015. 18

{[}97{]} Colin Raffel, Noam Shazeer, Adam Roberts, Katherine Lee, Sharan
Narang, Michael Matena, Yanqi Zhou, Wei Li, and Peter J Liu. Exploring
the limits of transfer learning with a unified text-to-text transformer.
The Journal of Machine Learning Research, 21(1):5485--5551, 2020. 23\\
{[}98{]} Jeff Rasley, Samyam Rajbhandari, Olatunji Ruwase, and Yuxiong
He. DeepSpeed: System optimizations enable training deep learning models
with over 100 billion parameters. In International Conference on
Knowledge Discovery and Data Mining (KDD), 2020. 24\\
{[}99{]} Ilya Loshchilov and Frank Hutter. Decoupled weight decay
regularization. In International Conference on Learning Representations
(ICLR), 2019. 24\\
{[}100{]} Mathlib Community. Mathport: A tool for porting Lean 3
projects to Lean 4. URL
https://github.com/leanprover-community/mathport.27\\
{[}101{]} OpenAI. ChatGPT plugins.
https://openai.com/blog/chatgpt-plugins, 2023. URL
https://openai.com/blog/chatgpt-plugins. 29\\
{[}102{]} Shunyu Yao, Dian Yu, Jeffrey Zhao, Izhak Shafran, Thomas L.
Griffiths, Yuan Cao, and Karthik Narasimhan. Tree of thoughts:
Deliberate problem solving with large language models. arXiv preprint
arXiv:2305.10601, 2023. 31, 33\\
{[}103{]} Erik Nijkamp, Bo Pang, Hiroaki Hayashi, Lifu Tu, Huan Wang,
Yingbo Zhou, Silvio Savarese, and Caiming Xiong. CodeGen: An open large
language model for code with multi-turn program synthesis. arXiv
preprint arXiv:2203.13474, 2022. 32\\
{[}104{]} Qinkai Zheng, Xiao Xia, Xu Zou, Yuxiao Dong, Shan Wang, Yufei
Xue, Zihan Wang, Lei Shen, Andi Wang, Yang Li, et al.~CodeGeeX: A
pre-trained model for code generation with multilingual evaluations on
HumanEval-X. arXiv preprint arXiv:2303.17568, 2023. 32\\
{[}105{]} Lili Yu, Dániel Simig, Colin Flaherty, Armen Aghajanyan, Luke
Zettlemoyer, and Mike Lewis. MegaByte: Predicting million-byte sequences
with multiscale transformers. arXiv preprint arXiv:2305.07185, 2023.
32\\
{[}106{]} Gautier Izacard and Édouard Grave. Leveraging passage
retrieval with generative models for open domain question answering. In
European Chapter of the Association for Computational Linguistics
(EACL), 2021. 33\\
{[}107{]} Yi Tay, Vinh Tran, Mostafa Dehghani, Jianmo Ni, Dara Bahri,
Harsh Mehta, Zhen Qin, Kai Hui, Zhe Zhao, Jai Gupta, et al.~Transformer
memory as a differentiable search index. In Neural Information
Processing Systems (NeurIPS), 2022. 34\\
{[}108{]} Michele Bevilacqua, Giuseppe Ottaviano, Patrick Lewis, Scott
Yih, Sebastian Riedel, and Fabio Petroni. Autoregressive search engines:
Generating substrings as document identifiers. 2022.\\
{[}109{]} Ronak Pradeep, Kai Hui, Jai Gupta, Adam D Lelkes, Honglei
Zhuang, Jimmy Lin, Donald Metzler, and Vinh Q Tran. How does generative
retrieval scale to millions of passages? arXiv preprint
arXiv:2305.11841, 2023. 34\\
{[}110{]} Norman Megill and David A Wheeler. Metamath: a computer
language for mathematical proofs. Lulu. com, 2019. 36

\section{Appendix}\label{appendix}

\section{A LeanDojo Technical Details
18}\label{a-leandojo-technical-details-18}

A.1 Extracting Premise Information from Lean's Elaborator \ldots.. 18\\
A.2 Reliable Interaction with Lean 18\\
A.3 Comparison with Existing Tools for Learning-Based Theorem Proving in
Lean . . . 19

\section{B LeanDojo Benchmark 20}\label{b-leandojo-benchmark-20}

B.1 Dataset Format 20\\
B.2 Datasheet 22\\
B.3 Data Hosting, Licensing, and Maintenance \ldots.. 23

\section{C Experiments 23}\label{c-experiments-23}

C.1 Details and Hyperparameters \ldots.. 23\\
C.2 The GPT-4 Baseline \ldots.. 24\\
C.3 Justifications for Not Comparing with Existing LLM-Based Provers
\ldots.. 26\\
C.4 Evaluation on MiniF2F and ProofNet 26

\section{D LeanDojo for Lean 4 27}\label{d-leandojo-for-lean-4-27}

\section{E ChatGPT Plugin for Theorem Proving
29}\label{e-chatgpt-plugin-for-theorem-proving-29}

\section{F Limitations and Future Work
32}\label{f-limitations-and-future-work-32}

\section{A LeanDojo Technical
Details}\label{a-leandojo-technical-details}

We provide more information on how LeanDojo extracts data from and
interacts with Lean. \(^{9}\) For further details, please check our
open-source implementation.

\section{A.1 Extracting Premise Information from Lean's
Elaborator}\label{a.1-extracting-premise-information-from-leans-elaborator}

``Premises'' in this paper belong to a category of Lean expressions
called ``constants.'' In Lean, definitions of constants are grouped into
nested, hierarchical namespaces. Therefore, each premise has a unique
fully-qualified name. For example, mod\_self in Fig. 2 is defined in the
namespace nat; therefore, its fully qualified name is nat.mod\_self.
However, it would be too verbose if premises had to be referred to using
full names. In practice, tactics often refer to premises using short
names such as mod\_self. In case multiple premises share the same short
name, Lean automatically infers the correct one from the context through
a process called ``name resolution''. LeanDojo is able to trace the
input/output of Lean's name resolution and thereby extract accurate
premise information for training the retriever.

Name resolution in Lean is implemented in a process called
``elaboration,'' which happens after parsing but before the parsed
expressions are checked by Lean's trusted kernel. Elaboration takes as
input user-entered expressions (called ``pre-expressions'') that are
concise, partially specified, and potentially ambiguous. It turns them
into complete expressions ready to be checked by the kernel. This is
realized by inferring not only full names but also missing types,
implicit arguments, overloading, type coercion, etc. Please refer to de
Moura et al.~{[}96{]} for details on Lean's elaboration process. In
LeanDojo, we modify Lean's internal implementation, intercepting the
elaborator to record its input/output:

\begin{itemize}
\tightlist
\item
  Pre-expression: The input to Lean's elaborator, including where
  premises are used in proofs.\\
\item
  Expression: The output of the elaborator, including the premise's full
  name and where it is defined.
\end{itemize}

Locations are spans in the source code, specified by the file name and
the row/column numbers of its start/end. Our modification takes the form
of a Git patch that LeanDojo can automatically apply to any version of
Lean 3 after March 24, 2022.

\section{A.2 Reliable Interaction with
Lean}\label{a.2-reliable-interaction-with-lean}

Polu et al.~{[}19{]} introduced lean-gym. To our knowledge, it is the
only mature, open-source tool before LeanDojo for interacting with Lean
programmatically. However, we found severe issues with lean-gym: About
21.1\% of the correct, human-written proofs are misjudged as incorrect,
leading to two problems: First, it underestimates the prover's
evaluation performance. Second, the results are too noisy as feedback
signals for reinforcement learning.

After carefully analyzing lean-gym's implementation, we identified the
root cause of the problem. When proving a theorem, the environment used
by lean-gym is subtly different from the original environment used by
humans. Specifically, lean-gym fails to handle namespaces correctly
(illustrated in Fig. A). As a result, name resolution fails unexpectedly
when checking correct proofs.

For example, Fig. A compares the correct environment and the environment
constructed by lean-gym. The theorem should be inside the namespace
``buffer''. However, in lean-gym, it merely opens the namespace. These
two scenarios are different when it comes to name resolution. Being
inside a namespace instructs Lean to favor constants defined in that
namespace, whereas opening a namespace does not have such an effect. In
this example, the short name ``read'' is ambiguous: We have
``monad\_reader.read'' defined in ``init/control/reader.lean'' and
``buffer.read'' defined in ``data/buffer.lean''. In the correct
environment, the ``read'' in ``unfold read'' resolves to
``buffer.read''. Whereas in lean-gym's environment, it incorrectly
resolved to ``monad\_reader.read''. Lean complains that ``read'' is not
an equational lemma, because it is referring to a wrong ``read''.
LeanDojo does not suffer from this kind of error since it uses a
different mechanism for constructing the environment. Specifically, it
wraps the interaction code as a Lean tactic, which is inserted into the
proof. Therefore, the environment is guaranteed to be correct.

We quantitatively compare lean-gym and LeanDojo on the number of proof
checking errors. In this study, we use Lean v3.42.1 paired with mathlib
version 6e5ca7d0097313e59f7533a42e3ea5197484c775 since they are
supported by both tools. We use LeanDojo to extract all tactic-style
proofs and enter them into both tools. These proofs are all correct, but
lean-gym failed on 21.1\% of them. In contrast, LeanDojo only failed on
1.4\%, and its failures are a subset of lean-gym's. We include this
study in our open-source repo and document example proofs from the
remaining 1.4\% to provide transparency on LeanDojo's limitations.
\(^{10}\)

\textit{[Image omitted: images/0478c53ab3d35c6b48b8a38bd2fa91a60e6db6c0dd500d2e72425309822b129d.jpg]}

text\_image

import data.buffer universe u

namespace buffer

theorem my\_read\_eq\_read' \{a : Type u\} {[}inhabited a{]} (b : buffer
a) (i : nat) (h : i \textless{} b.size) : read b \textless i,
h\textgreater{} = read' b i := begin

cases b, unfold read, unfold read', simp {[}array.read\_eq\_read'{]} end

end buffer

Correct environment

import data.buffer universe u

open buffer

theorem my\_read\_eq\_read' \{a : Type u\} {[}inhabited a{]} (b : buffer
a) (i : nat) (h : i \textless{} b.size) : read b \textless i,
h\textgreater{} = read' b i := begin

cases b, unfold read, unfold read', simp {[}array.read\_eq\_read'{]} end

lean-gym's environment

ERROR: unfold tactic failed, \texttt{read} does not have equational
lemmas nor is a projection

Figure A: An example of correct proofs misjudged as incorrect by
lean-gym, adapted from the theorem read\_eq\_read' in
``data/buffer.lean'' of Lean's standard library. The error message is
because lean-gym failed to resolve the short name ``read'' to the
correct fully-qualified name. The Lean code in this figure is only for
illustrative purposes. It does not reflect the implementation technique
used by lean-gym to construct the environment. Instead of generating
actual Lean code, lean-gym uses Lean's metaprogramming APIs to construct
the environment.

\section{A.3 Comparison with Existing Tools for Learning-Based Theorem
Proving in
Lean}\label{a.3-comparison-with-existing-tools-for-learning-based-theorem-proving-in-lean}

To our knowledge, LeanStep {[}16{]} \(^{11}\) and lean-gym {[}19{]} are
the only published tools for learning-based theorem proving in Lean.
There are a few unpublished prototypes, such as repl,
lean-client-python, and lean-gym for Lean 4, none of which is mature
enough or is under active development. Therefore, we only compare
LeanDojo with LeanStep and lean-gym (summarized in Table A).

Functionality. LeanDojo supports both data extraction and interacting
with Lean programmatically. In contrast, LeanStep is only for data
extraction, and lean-gym is only for interacting with Lean. They are not
actively maintained, so they do not support recent versions of mathlib
(tested on August 11, 2023, using mathlib commit
19c869efa56bbb8b500f2724c0b77261edbfa28c). Also, neither of them support
Lean 4 (Appendix D). LeanDojo fully supports recent mathlib and Lean 4.
Furthermore, LeanStep cannot extract premise information and is not
applicable to repos other than mathlib. Last, LeanDojo comes with
comprehensive documentation and unit tests, whereas other tools barely
have any.

Implementation details. LeanStep and LeanDojo use different mechanisms
to extract ASTs and proof trees. LeanStep implements an ad-hoc parser in
Python for parsing Lean code into ASTs. It also intercepts Lean's tactic
system to insert logging code. Then the logs are used to reconstruct
proof trees. This implementation is brittle and does not work for the
current versions of Lean/mathlib. In contrast, LeanDojo relies on Lean's
built-in mechanisms for exporting ASTs and proof states (lean --ast
--tsast --tspp), which works robustly for recent Lean/mathlib. This
mechanism was developed after LeanStep.

Regarding interaction with Lean, both lean-gym and LeanDojo rely on
Lean's metaprogramming APIs, and LeanDojo partially builds upon
lean-gym's code. However, lean-gym has a critical issue in that it
misjudges many correct proofs as incorrect (Appendix A.2). The main
reason is that lean-gym fails to distinguish two subtly different cases
when constructing the proof environment: (1) opening a namespace; (2)
being inside a namespace. LeanDojo does not suffer from this issue.
Instead of operating as a standalone program in the IO monad, it wraps
the interaction code into a special tactic, which is inserted into the
correct location in the proof. Therefore, the interaction code is
guaranteed to run in the same environment as the original human-written
proof.

LeanStep {[}16{]}

lean-gym {[}19{]}

LeanDojo (ours)

Data extraction

Premise information

✗

N/A

√

Lean 4 support

✗

N/A

√

Recent mathlib

✗

N/A

√

Repos other than mathlib

✗

N/A

√

Interaction

Estimated errors

N/A

21.1\%

1.4\%

Lean 4 support

N/A

✗

√

Recent mathlib

N/A

✗

√

Repos other than mathlib

N/A

√

√

Documentation \& unit tests

✗

✗

√

Table A: Comparing LeanDojo with existing tools for data extraction and
interaction with Lean.

\section{B LeanDojo Benchmark}\label{b-leandojo-benchmark}

\section{B.1 Dataset Format}\label{b.1-dataset-format}

We describe the data format of LeanDojo Benchmark, which has the
following directory structure:

\begin{Shaded}
\begin{Highlighting}[]
\NormalTok{corpus.jsonl......All premises defined in mathlib and Lean\textquotesingle{}s standard library}
\NormalTok{metadata.json......Metadata}
\NormalTok{licenses}
\NormalTok{lean......Attribution to Lean\textquotesingle{}s Apache 2.0 license}
\NormalTok{mathlib......Attribution to mathlib\textquotesingle{}s Apache 2.0 license}
\NormalTok{README.md......Statement that LeanDojo Benchmark is released under CC BY 2.0}
\NormalTok{random......Theorems/proofs of the random split}
\NormalTok{train.json......94,734 theorems}
\NormalTok{val.json......2,000 theorems}
\NormalTok{test.json......2,000 theorems}
\NormalTok{novel\_premises......Theorems/proofs of the novel\_premises split}
\NormalTok{train.json......94,734 theorems}
\NormalTok{val.json......2,000 theorems}
\NormalTok{test.json......2,000 theorems }
\end{Highlighting}
\end{Shaded}

Premise Definitions. corpus.jsonl contains the definition of premises.
It has 3,280 lines. Each line is in JSON format and corresponds to a
Lean file. Below is an example for ``init/control/functor.lean'', which
directly imports three other files: ``init/core.lean'',
``init/function.lean'', and ``init/meta/name.lean''. It defines two
constants that can be used as premises: ``functor'' and
``functor.map\_const\_rev''. For each premise, we have access to its
full name, the source code, and its start/end location within the file.

\begin{Shaded}
\begin{Highlighting}[]
\CommentTok{"path"}\NormalTok{: }\StringTok{"\_target/deps/lean/library/init/control/functor.lean"}\NormalTok{,}
\CommentTok{"imports"}\NormalTok{: [}
    \StringTok{"\_target/deps/lean/library/init/core.lean"}\NormalTok{,}
    \StringTok{"\_target/deps/lean/library/init/function.lean"}\NormalTok{,}
    \StringTok{"\_target/deps/lean/library/init/meta/name.lean"}
\NormalTok{],}
\CommentTok{"premises"}\NormalTok{: [ }
\end{Highlighting}
\end{Shaded}

\begin{Shaded}
\begin{Highlighting}[]
\FunctionTok{\{}
    \DataTypeTok{"full\_name"}\FunctionTok{:} \StringTok{"functor"}\FunctionTok{,}
    \DataTypeTok{"code"}\FunctionTok{:} \StringTok{"class functor (f : Type u → Type v) : Type (max (u+1) v) := }\CharTok{\textbackslash{}n}\StringTok{(map : Π \{α β : Type u\}, (α → β) → f α → f β) }\CharTok{\textbackslash{}n}\StringTok{(map\_const : Π \{α β : Type u\}, α → f β → f α := λ α β, map o const β)"}\FunctionTok{,}
    \DataTypeTok{"start"}\FunctionTok{:} \OtherTok{[}\DecValTok{11}\OtherTok{,} \DecValTok{1}\OtherTok{]}\FunctionTok{,}
    \DataTypeTok{"end"}\FunctionTok{:} \OtherTok{[}\DecValTok{13}\OtherTok{,} \DecValTok{70}\OtherTok{]}\FunctionTok{,}
    \DataTypeTok{"kind"}\FunctionTok{:} \StringTok{"class"}
\FunctionTok{\}}\ErrorTok{,}
\FunctionTok{\{}
    \DataTypeTok{"full\_name"}\FunctionTok{:} \StringTok{"functor.map\_const\_rev"}\FunctionTok{,}
    \DataTypeTok{"code"}\FunctionTok{:} \StringTok{"@[reducible] def functor.map\_const\_rev \{f : Type u → Type v\} [functor f] \{α β : Type u\} : f β → α → f α := }\CharTok{\textbackslash{}n}\StringTok{λ a b, b \textless{}$ a"}\FunctionTok{,}
    \DataTypeTok{"start"}\FunctionTok{:} \OtherTok{[}\DecValTok{18}\OtherTok{,} \DecValTok{1}\OtherTok{]}\FunctionTok{,}
    \DataTypeTok{"end"}\FunctionTok{:} \OtherTok{[}\DecValTok{19}\OtherTok{,} \DecValTok{14}\OtherTok{]}\FunctionTok{,}
    \DataTypeTok{"kind"}\FunctionTok{:} \StringTok{"definition"}
\FunctionTok{\}} 
\end{Highlighting}
\end{Shaded}

Theorems and Tactics. Theorems in LeanDojo Benchmark are split into
training/validation/testing using two different strategies (Sec. 4).
They are formatted in JSON, and below is an example corresponding to the
theorem ``real.angle.to\_real(pi\_div\_two''. LeanDojo has recorded two
tactics: ``split'' and ``linarith {[}pi\_pos{]}''. For each tactic, we
have the proof states before/after it. The ``linarith {[}pi\_pos{]}''
tactic illustrates how premises are recorded: They are annotated using
HTML-like strings such as ``linarith {[}pi\_pos{]}'', followed by a
``provenance list''. Each element in the list corresponds to a premise
in the tactic.

\begin{Shaded}
\begin{Highlighting}[]
\StringTok{"url"}\ErrorTok{:} \StringTok{"https://github.com/leanprover{-}community/mathlib"}\ErrorTok{,}
\StringTok{"commit"}\ErrorTok{:} \StringTok{"19c869efa56bbb8b500f2724c0b77261edbfa28c"}\ErrorTok{,}
\StringTok{"file\_path"}\ErrorTok{:} \StringTok{"src/analysis/special\_functions/trigonometric/angle.lean"}\ErrorTok{,}
\StringTok{"full\_name"}\ErrorTok{:} \StringTok{"real.angle.to\_real(pi\_div\_two"}\ErrorTok{,}
\StringTok{"start"}\ErrorTok{:} \OtherTok{[}\DecValTok{512}\OtherTok{,} \DecValTok{9}\OtherTok{]}\ErrorTok{,}
\StringTok{"end"}\ErrorTok{:} \OtherTok{[}\DecValTok{513}\OtherTok{,} \DecValTok{56}\OtherTok{]}\ErrorTok{,}
\StringTok{"traced\_tactics"}\ErrorTok{:} \OtherTok{[}
    \FunctionTok{\{}
    \DataTypeTok{"tactic"}\FunctionTok{:} \StringTok{"split"}\FunctionTok{,}
    \DataTypeTok{"annotated\_tactic"}\FunctionTok{:} \OtherTok{[}\StringTok{"split"}\OtherTok{,} \OtherTok{[]]}\FunctionTok{,}
    \DataTypeTok{"state\_before"}\FunctionTok{:} \StringTok{"|{-}π \textless{} π / 2 ∧ π / 2 ≤ π"}\FunctionTok{,}
    \DataTypeTok{"state\_after"}\FunctionTok{:} \StringTok{"2 goals}\CharTok{\textbackslash{}n}\StringTok{|{-}π \textless{} π / 2}\CharTok{\textbackslash{}n\textbackslash{}n}\StringTok{|{-}π / 2 ≤ π"}
    \FunctionTok{\}}\OtherTok{,}
    \FunctionTok{\{}
    \DataTypeTok{"tactic"}\FunctionTok{:} \StringTok{"linarith [pi\_pos]"}\FunctionTok{,}
    \DataTypeTok{"annotated\_tactic"}\FunctionTok{:} \OtherTok{[}
    \StringTok{"linarith [\textless{}a\textgreater{}pi\_pos\textless{}/a\textgreater{}]"}\OtherTok{,} \OtherTok{[}
    \FunctionTok{\{}
    \DataTypeTok{"full\_name"}\FunctionTok{:} \StringTok{"real(pi\_pos"}\FunctionTok{,}
    \DataTypeTok{"def\_path"}\FunctionTok{:} \StringTok{"src/analysis/special\_functions/trigonometric/basic.lean"}\FunctionTok{,}
    \DataTypeTok{"def\_pos"}\FunctionTok{:} \OtherTok{[}\DecValTok{122}\OtherTok{,} \DecValTok{7}\OtherTok{]}\FunctionTok{,}
    \FunctionTok{\}}
    \OtherTok{]}
    \OtherTok{]}\FunctionTok{,}
    \DataTypeTok{"state\_before"}\FunctionTok{:} \StringTok{"|{-}π \textless{} π / 2"}\FunctionTok{,}
    \DataTypeTok{"state\_after"}\FunctionTok{:} \StringTok{"no goals"}
    \FunctionTok{\}}
\OtherTok{]} 
\end{Highlighting}
\end{Shaded}

Not all theorems have tactic-style proofs. For those without
tactic-style proofs, concatenating the tactics does not lead to a
complete proof of the original theorem. However, this is not an issue
when using the data for theorem proving evaluation or for training
tactic generators.

\section{B.2 Datasheet}\label{b.2-datasheet}

We present a datasheet {[}90{]} for documentation and responsible usage
of LeanDojo Benchmark.

\section{Motivation.}\label{motivation.}

\begin{itemize}
\tightlist
\item
  For what purpose was the dataset created? It was created as a
  benchmark for learning-based theorem proving in Lean.\\
  • Who created the dataset (e.g., which team, research group) and on
  behalf of which entity (e.g., company, institution, organization)? It
  was created by the authors of this paper.\\
  • Who funded the creation of the dataset? See the acknowledgments in
  Sec. 7.
\end{itemize}

\section{Composition.}\label{composition.}

\begin{itemize}
\tightlist
\item
  What do the instances that comprise the dataset represent (e.g.,
  documents, photos, people, countries)? The dataset consists of formal
  definitions, theorems, and proofs written in Lean {[}1{]}.\\
\item
  How many instances are there in total (of each type, if appropriate)?
  The dataset has 98,734 theorems and their proofs, as well as 130,262
  premises defined in 3,384 files.\\
\item
  Does the dataset contain all possible instances or is it a sample (not
  necessarily random) of instances from a larger set? The dataset
  contains all theorems/proofs that LeanDojo can extract from the commit
  19c869efa56bbb8b500f2724c0b77261edbfa28c of mathlib released on
  October 11, 2023.\\
\item
  What data does each instance consist of? Theorems/proofs in the
  dataset are Lean code written by programmers and mathematicians.\\
\item
  Are relationships between individual instances made explicit?
  Definitions in the dataset are linked to proofs using them as
  premises.\\
\item
  Are there recommended data splits? Yes, we recommend two data splits:
  random and novel\_premises. Please see Sec. 4 for details.\\
\item
  Are there any errors, sources of noise, or redundancies in the
  dataset? ASTs extracted by LeanDojo contain a small number of errors
  due to potential flaws in Lean's AST exporting mechanism. However,
  they do not have a tangible impact on our work.\\
\item
  Is the dataset self-contained, or does it link to or otherwise rely on
  external resources (e.g., websites, tweets, other datasets)? The
  dataset is self-contained.\\
\item
  Does the dataset contain data that might be considered confidential
  (e.g., data that is protected by legal privilege or by doctor-patient
  confidentiality, data that includes the content of individuals'
  non-public communications)? No.\\
\item
  Does the dataset contain data that, if viewed directly, might be
  offensive, insulting, threatening, or might otherwise cause anxiety?
  No.
\end{itemize}

\section{Collection Process.}\label{collection-process.}

\begin{itemize}
\item
  How was the data associated with each instance acquired? The data is
  directly observable by opening mathlib in VS Code with the Lean
  plugin. However, we had to instrument Lean to export the data
  programmatically.\\
\item
  What mechanisms or procedures were used to collect the data (e.g.,
  hardware apparatuses or sensors, manual human curation, software
  programs, software APIs)? The data was generated by building a Lean
  repo using our modified Lean and postprocessing the exported data.
\item
  Who was involved in the data collection process (e.g., students, crowd
  workers, contractors), and how were they compensated (e.g., how much
  were crowd workers paid)? No manual effort was involved in the data
  collection process.\\
  • Over what timeframe was the data collected? The final version of the
  dataset was generated in October 2023.
\end{itemize}

\section{Uses.}\label{uses.}

\begin{itemize}
\tightlist
\item
  Has the dataset been used for any tasks already? We have used the
  dataset for training and evaluating machine learning models on the
  tasks of premise selection and theorem proving.\\
\item
  Is there a repository that links to any or all papers or systems that
  use the dataset? Yes, https://leandojo.org.
\end{itemize}

\section{Distribution.}\label{distribution.}

\begin{itemize}
\tightlist
\item
  Will the dataset be distributed to third parties outside of the entity
  (e.g., company, institution, organization) on behalf of which the
  dataset was created? Yes, the dataset is publicly available on the
  Internet.\\
\item
  How will the dataset be distributed (e.g., tarball on website, API,
  GitHub)? The dataset can be downloaded as a tarball.\\
\item
  Will the dataset be distributed under a copyright or other
  intellectual property (IP) license, and/or under applicable terms of
  use (ToU)? The dataset is distributed under CC BY 2.0. The data
  generation code is distributed under the MIT license. The dataset was
  extracted from mathlib, which depends on lean. Both of them are
  distributed under the Apache 2.0 license. We include their licenses in
  the dataset as attribution (Appendix B.1).\\
\item
  Have any third parties imposed IP-based or other restrictions on the
  data associated with the instances? No.\\
\item
  Do any export controls or other regulatory restrictions apply to the
  dataset or to individual instances? No.
\end{itemize}

\section{Maintenance.}\label{maintenance.}

• Who will be supporting/hosting/maintaining the dataset? The authors of
this paper.\\
- How can the owner/curator/manager of the dataset be contacted (e.g.,
email address)? Please contact Kaiyu Yang at kaiyuy@caltech.edu.\\
- Is there an erratum? No.\\
- Will the dataset be updated (e.g., to correct labeling errors, add new
instances, delete instances)? Please check https://leandojo.org for any
update.\\
- If others want to extend/augment/build on/contribute to the dataset,
is there a mechanism for them to do so? Yes, they can use our data
generation code, which is publicly available.

\section{B.3 Data Hosting, Licensing, and
Maintenance}\label{b.3-data-hosting-licensing-and-maintenance}

LeanDojo Benchmark is distributed under the CC BY 2.0 license. The data
is hosted on zenodo.org (a long-term data repository operated by CERN).
The LeanDojo tool for data extraction and interaction with Lean is
released at https://github.com/lean-dojo/LeanDojo under the MIT license.
Our model checkpoints are hosted on Hugging Face Hub. LeanDojo's
documentation is hosted on Read the Docs at
https://leandojo.readthedocs.io. LeanDojo's website
(https://leandojo.org) is the entry point for everything related to it,
including any future updates or maintenance.

\section{C Experiments}\label{c-experiments}

\section{C.1 Details and
Hyperparameters}\label{c.1-details-and-hyperparameters}

The premise retriever and tactic generator in ReProver are initialized
by the google/byt5-small checkpoint on Hugging Face. It is a T5-like
{[}97{]} encoder-decoder Transformer that operates directly

on UTF-8 bytes without tokenization. We choose ByT5 {[}44{]} instead of
T5 because Lean code makes extensive use of Unicode math symbols, which
may cause problems to T5's pretrained tokenizer. The retriever uses the
encoder only, whereas the generator uses both the encoder and the
decoder.

In training, we use one NVIDIA A100 GPU with 80GB of memory. The code is
implemented in PyTorch and PyTorch Lightning, with bfloat16 mixed
precision and DeepSpeed ZeRO Stage 2 {[}98{]}. Both the retriever and
the generator are optimized using AdamW {[}99{]} with a batch size of 8.
In the first 2,000 steps, the learning rate warms up linearly from 0 to
the maximum value. Then it decays to 0 following a cosine schedule. The
maximum learning rate is \(10^{-4}\) for the retriever and
\(5 \times 10^{-4}\) for the generator. When training the retriever, we
sample 3 negative premises for each example, including 1 in-file
negative premise. When training the generator, we apply dropout to
retrieved premises with a dropout rate of 0.5. Then, we truncate the
generator's input to 2,300 tokens.

During evaluation, the tactic generator is combined with best-first
search to find proofs. At each search step, it produces 64 tactic
candidates using beam search. Each tactic is associated with a
log-likelihood score. In best-first search, we prioritize the states by
the sum of log-likelihoods of tactics leading to that state.

\section{C.2 The GPT-4 Baseline}\label{c.2-the-gpt-4-baseline}

Now we describe the GPT-4 {[}27{]} baseline in Sec. 6. Similar to
ReProver, it is a tactic generator combined with best-first search.
However, the tactic generator is based on GPT-4's capability to follow
instructions in zero-shot. Specifically, given a proof state, we use the
following prompt to instruct GPT-4 to produce a list of tactics, each
paired with a confidence score:

\section{Prompt Template:}\label{prompt-template}

You are an expert in Lean3 theorem proofs. We are trying to solve the
Lean3 theorem `THEOREM\_FULL\_NAME' from the mathlib file `FILE\_PATH'.
The current tactic state is: `TACTIC\_STATE'. Suggest exactly 35 unique
tactics to progress in solving `THEOREM\_FULL\_NAME', along with their
confidence levels as a float between 0 and 1. Rank them in order of
effectiveness. Present the tactics and their confidence levels as
comma-separated tuples in this format: \#(tactic\_\{1\},
confidence\_\{1\})\#, \#(tactic\_\{2\}, confidence\_\{2\})\#, \ldots,
\#(tactic\_\{35\}, confidence\_\{35\})\#.

We adapted the prompt to a particular theorem and state by substituting
the variables with the appropriate values. Given the inherent
variability in GPT-4's outputs, we requested 35 and filtered out invalid
ones. We used a token length limit of 1,024 and kept all other API
parameters at their default values. Below are a few example prompts and
GPT-4's responses:

\section{Example Prompt 1:}\label{example-prompt-1}

You are an expert in Lean3 theorem proofs. We are trying to solve the
Lean3 theorem `polynomial.chebyshev.aeval\_U' from the mathlib file
`mathlib/src/analysis/special\_functions/trigonometric/chebyshev.lean'.
The current tactic state is: `R : Type u\_1, A : Type u\_2, \_inst\_1 :
comm\_ring R, \_inst\_2 : comm\_ring A, \_inst\_3 : algebra R A, x : A,
n : N ⊢ ↑(aeval x) (chebyshev.U R n) = eval x (chebyshev.U A n)'.
Suggest exactly 35 unique tactics to progress in solving
`polynomial.chebyshev.aeval\_U', along with their confidence levels as a
float between 0 and 1. Rank them in order of effectiveness. Present the
tactics and their confidence levels as comma-separated tuples in this
format: \#(tactic\_\{1\}, confidence\_\{1\})\#, \#(tactic\_\{2\},
confidence\_\{2\})\#, \ldots, \#(tactic\_\{35\}, confidence\_\{35\})\#.

\section{Example Response 1:}\label{example-response-1}

\# (rw chebyshev.U, 0.9)\#, \#(simp only {[}chebyshev.U{]}, 0.9)\#,
\#(simp\_rw chebyshev.U, 0.8)\#, \#(apply aeval\_add, 0.75)\#, \#(apply
aeval\_mul, 0.75)\#, \#(apply aeval\_sub, 0.75)\#, \#(apply aeval\_X,
0.7)\#, \#(apply aeval\_C, 0.7)\#, \#(apply aeval\_pow, 0.7)\#, \#(rw
\u2190aeval\_pow, 0.65)\#, \#(rw 2190aeval\_mul, 0.65)\#, \#(rw
\u2190aeval\_add, 0.65)\#, \#(rw \u2190aeval\_sub, 0.65)\#, \#(apply
algebra\_map\_eq, 0.6)\#, \#(intro aeval\_algebra\_map\_eq, 0.6)\#,
\#(apply eval\_map, 0.55)\#, \#(rw eval\_map, 0.55)\#, \#(apply
aeval\_eq\_eval, 0.55)\#, \#(rw\_for\_target aeval\_eq\_eval, 0.5)\#,
\#(apply aeval\_eq\_eval', 0.5)\#, \#(rw\_for\_target aeval\_eq\_eval',
0.5)\#, \#(apply aeval\_commutes, 0.45)\#, \#(rw aeval\_commutes,
0.45)\#, \#(apply aeval\_eq\_of\_is\_scalar\_tower, 0.4)\#, \#(rw
is\_scalar\_tower.aeval\_apply, 0.4)\#, \#(try\_for 100 exact\_mod\_cast
chebyshev.U\_eq\_ref1, 0.35)\#, \#(simp * at *, 0.3)\#, \#(simp
int.coe\_nat\_ring\_hom, 0.25)\#, \#(apply chebyshev.U\_is\_polynomial,
0.2)\#, \#(rw\_algebra\_map\_eq\_of\_ring\_hom, 0.15)\#, \#(unfold
chebyshev.U, 0.1)\#, \#(congr, 0.05)\#, \#(change (chebyshev.U R n).eval
x = chebyshev.U A n, 0.05)\#, \#(change (\u21d1(aeval x) (chebyshev.U R
n) = chebyshev.U A n), 0.05)\#, \#(apply\_fun aeval x, 0.05)\#,
\#(apply\_fun eval x, 0.05)\#, \#(simp context\_md := reducible {[}{]},
0.01)\#

\section{Example Prompt 2:}\label{example-prompt-2}

You are an expert in Lean3 theorem proofs. We are trying to solve the
Lean3 theorem `euler\_sine.integral\_cos\_pow\_pos' from the math-lib
file
`mathlib/src/analysis/special\_functions/trigonometric/euler\_sine\_prod.lean'.
The current tactic state is: `n : N ⊢ 0 \textless{} ∫ (x : R) in 0..π /
2, cos x \^{} n'. Suggest exactly 35 unique tactics to progress in
solving `euler\_sine.integral\_cos\_pow\_pos', along with their
confidence levels as a float between 0 and 1. Rank them in order of
effectiveness. Present the tactics and their confidence levels as
comma-separated tuples in this format: \#(tactic\_1, confidence\_1)\#,
\#(tactic\_2, confidence\_2)\#, \ldots, \#(tactic\_35,
confidence\_35)\#.

\section{Example Response 2:}\label{example-response-2}

\# (exact integral\_pos\_of\_antimono\_increasing\_on
cos.pow\_antimono\_increasing\_on\_cos\_even\_nnreal, 0.99)\#, \#(apply
integral\_pos\_of\_antimono\_increasing\_on, 0.95)\#,
\#(interval\_integrable\_integral\_pos, 0.90)\#, \#(exact
integral\_pos\_pre, 0.85)\#, \#(interval\_integ\_timedout, 0.80)\#,
\#(interval\_integrable\_pos, 0.75)\#, \#(exact
integral\_pos\_of\_nonneg\_dint\_val, 0.70)\#, \#(exact
interval\_integrable\_pos\_left, 0.65)\#, \#(exact
interval\_integrable\_pos\_right, 0.60)\#, \#(exact
integral\_pos\_of\_nonneg\_f, 0.55)\#, \#(rw integral\_on\_interval,
0.50)\#, \#(rw interval\_integral\_pos\_of\_supporting, 0.45)\#, \#(rw
integral\_pos\_of\_interval\_integrable\_pos, 0.40)\#, \#(apply
integral\_pos\_of\_antimono\_increasing\_on\_cos\_even\_nnreal, 0.35)\#,
\#(norm\_cast apply\_pos, 0.30)\#, \#(apply\_pos\_of\_nonneg\_pos,
0.25)\#, \#(rw interval\_integral\_pos\_of\_antimono, 0.20)\#, \#(use
cos\_slow\_mono\_on\_pos, 0.15)\#, \#(apply cos\_slow\_mono\_on\_pos,
0.10)\#, \#(exact integral\_pos\_pre\_cos\_pow\_even, 0.09)\#,
\#(norm\_cast apply\_nonneg\_pos, 0.08)\#, \#(rw
integral\_pos\_of\_monotone\_integrable, 0.07)\#, \#(use
cos\_pos\_constant\_sequence, 0.06)\#, \#(apply
cos\_pos\_constant\_sequence, 0.05)\#, \#(use
integral\_pos\_of\_ascending\_sequence, 0.04)\#, \#(apply
integral\_pos\_of\_ascending\_sequence, 0.03)\#,
\#(interval\_integral\_pos\_of\_antimono\_increasing\_on\_pos, 0.02)\#,
\#(norm\_cast apply\_pos\_of\_nonzero\_pos, 0.01)\#,
\#(apply\_pos\_of\_nonzero\_pos, 0.005)\#, \#(rw
integral\_pos\_of\_nonneg\_f\_monotone, 0.001)\#

Data contamination is possible. Our GPT-4 experiments were performed in
2023, but many theorems and proofs in the dataset have been publicly
available on GitHub before GPT-4's data cutoff date (September 2021).

\section{C.3 Justifications for Not Comparing with Existing LLM-Based
Provers}\label{c.3-justifications-for-not-comparing-with-existing-llm-based-provers}

In Table 2, we do not empirically compare ReProver with any existing
LLM-based prover. Unfortunately, such a comparison is infeasible.
Provers targeting different proof assistants are generally not
comparable, so we focus the discussion on the three existing provers in
Lean {[}16, 17, 19{]}. Most importantly, they are impossible to
reproduce with reasonable effort, due to private code and pretraining
data. Therefore, the only potential comparison is to evaluate ReProver
under their experimental settings and compare with the numbers reported
in their papers. However, that is also impractical for numerous reasons:

\begin{itemize}
\tightlist
\item
  The data is different. All existing methods used an outdated version
  of mathlib more than two years ago. We cannot use LeanDojo to extract
  data from this version. As mentioned in Sec. 4, LeanDojo only supports
  repos released after March 24, 2022. Also, we cannot use their dataset
  directly, since it does not contain premise information required by
  ReProver.\\
\item
  Lample et al.~{[}17{]} trained on a synthetic dataset named Equations,
  which is not publicly available.\\
\item
  All existing methods co-train the tactic generator on auxiliary tasks
  from the PACT dataset {[}16{]}. Co-training increases the data/compute
  requirements by an order of magnitude, which cannot be afforded by us
  (or probably most academic labs). All existing methods were developed
  by researchers in the industry.\\
\item
  Polu et al.~{[}19{]} and Lample et al.~{[}17{]} further finetuned
  their models on new proofs collected through online interaction with
  Lean, whereas our method is only trained on human-written proofs.\\
\item
  The tool for interacting with Lean may impact the performance. Han et
  al.~{[}16{]} and Polu et al.~{[}19{]} used lean-gym, which has severe
  limitations (Appendix A.2). Lample et al.~{[}17{]} developed their own
  private tool, which is not publicly available.
\end{itemize}

Most of these difficulties are due to the private nature of existing
methods. By releasing our code and models, we take a major step in
establishing accessible baselines for future work to build upon.

\section{C.4 Evaluation on MiniF2F and
ProofNet}\label{c.4-evaluation-on-minif2f-and-proofnet}

We evaluate our ReProver model on MiniF2F {[}28{]} and ProofNet {[}29{]}
(Sec. 6) to test its capability in proving theorems outside its training
data distribution. We use the same hyperparameters and evaluation setup
as the previous experiments (Appendix C.1).

MiniF2F. We use the commit 5271ddec788677c815cf818a06f368ef6498a106 of
Meta's version of MiniF2F {[}17{]}. ReProver achieves a Pass@1 of 26.5\%
on the test set, which is competitive with state-of-the-art methods
without reinforcement learning (25.9\% in Polu et al.~{[}19{]}).
Moreover, ReProver can prove 33 theorems that currently do not have Lean
proofs (examples in Fig. B). For the complete list of 33 new proofs,
please see our pull request to MiniF2F.

There are caveats about quantitatively comparing ReProver with existing
methods on MiniF2F. Many difficulties in Appendix C.3 still apply, e.g.,
different tools for interacting with Lean may impact the performance.
Also, MiniF2F is a test-only dataset without training theorems, and
existing methods focus on reinforcement learning (RL) to learn from
proofs collected via online interaction with the proof assistant {[}17,
19{]}. In contrast, ReProver is trained via supervised learning on a
static dataset, so we only compare with the non-RL baseline in existing
methods (Polu et al.~{[}19{]} achieves a Pass@1 of 25.9\% without RL and
29.6\% with RL). Furthermore, we do not compare with Lample et
al.~{[}17{]} due to differences in the evaluation metric. They use
Pass@64, which requires running the prover on each theorem 64 times. We
use Pass@1, and it already takes one day for a single evaluation on
MiniF2F's test set. Therefore, evaluating Pass@64 would be too
computationally expensive for the resources we have access to. Finally,
MiniF2F is available in multiple proof assistants {[}18, 69, 70{]}.
Results across different proof assistants are not comparable, so we only
compare with existing work in Lean.

ProofNet. We use the commit e8645aa830ce17c33a8b8482a8195f0f97d6a74a of
ProofNet. Re-Prover can prove 48 out of 349 theorems, achieving a Pass@1
of 13.8\%, which is the first reported

\begin{Shaded}
\begin{Highlighting}[]
\NormalTok{theorem mathd\_numbertheory\_237 }\OperatorTok{:}
\NormalTok{(∑ k }\ControlFlowTok{in}\NormalTok{ (finet}\AttributeTok{.range} \DecValTok{101}\NormalTok{), k) }\OperatorTok{\%} \DecValTok{6} \OperatorTok{=} \DecValTok{4} \OperatorTok{:=}
\ControlFlowTok{begin}
\NormalTok{rw }\KeywordTok{[}\NormalTok{finet}\AttributeTok{.sum\_range\_succ}\VerbatimStringTok{\textquotesingle{}],}
\VerbatimStringTok{norm\_num [finet.sum\_range\_succ],}
\VerbatimStringTok{end}

\VerbatimStringTok{theorem mathd\_numbertheory\_175 :}
\VerbatimStringTok{(2\^{}2010) \% 10 = 4 :=}
\VerbatimStringTok{begin}
\VerbatimStringTok{norm\_num [pow\_succ],}
\VerbatimStringTok{end}

\VerbatimStringTok{theorem mathd\_numbertheory\_293}
\VerbatimStringTok{(n : N)}
\VerbatimStringTok{(h₀ : n ≤ 9)}
\VerbatimStringTok{(h₁ : 11|20 * 100 + 10 * n + 7) :}
\VerbatimStringTok{n = 5 :=}
\VerbatimStringTok{begin}
\VerbatimStringTok{contrapose! h₁,}
\VerbatimStringTok{norm\_num [h₁],}
\VerbatimStringTok{dec\_trivial!,}
\VerbatimStringTok{end}

\VerbatimStringTok{theorem mathd\_algebra\_616}
\VerbatimStringTok{(f g : R → R)}
\VerbatimStringTok{(h₀ : ∀ x, f x = x\^{}3 + 2 * x + 1)}
\VerbatimStringTok{(h₁ : ∀ x, g x = x {-} 1) :}
\VerbatimStringTok{f (g 1) = 1 :=}
\VerbatimStringTok{begin}
\VerbatimStringTok{simp only [h₀, h₁, pow\_one],}
\VerbatimStringTok{ring,}
\VerbatimStringTok{end }
\end{Highlighting}
\end{Shaded}

Figure B: Examples of new proofs discovered by ReProver on MiniF2F
{[}28{]}.

theorem proving result on ProofNet. Moreover, 39 out of the 48 proved
theorems do not have existing Lean proofs (examples in Fig. C), and 3 of
them can only be proved with the help of premise retrieval (Fig. D). We
have contributed the 39 new proofs to ProofNet, which helped them reveal
and fix problems in the formalization of 7 theorems (details in our pull
request).

\section{D LeanDojo for Lean 4}\label{d-leandojo-for-lean-4}

Lean 3 and Lean 4 are two incompatible major versions of Lean, \(^{12}\)
and both are widely used. Lean 3 was the latest stable version until
recently (June 2023). Also, Lean 3 and Lean 4 have separate versions of
mathlib. The Lean/mathlib community has recently finished porting
theorems and proofs from mathlib 3 to mathlib 4 {[}100{]}. Therefore,
Lean 3 will gradually become deprecated, and future Lean projects will
be using Lean 4. Therefore, it is important for LeanDojo to support Lean
4.

Since Lean 4 is relatively new, we are not aware of any existing work on
learning-based theorem proving in Lean 4. Furthermore, no existing tool
is available for extracting data from Lean 4. LeanDojo fills in this gap
and fully supports Lean 4. Given any repo in Lean 4, LeanDojo can
extract

\begin{Shaded}
\begin{Highlighting}[]
\NormalTok{theorem exercise\_2\_3\_2 \{G }\OperatorTok{:} \DataTypeTok{Type}\OperatorTok{*}\NormalTok{\} [group G] (a b }\OperatorTok{:}\NormalTok{ G) }\OperatorTok{:}
\NormalTok{    ∃ g }\OperatorTok{:}\NormalTok{ G, b }\OperatorTok{*}\NormalTok{ a }\OperatorTok{=}\NormalTok{ g }\OperatorTok{*}\NormalTok{ a }\OperatorTok{*}\NormalTok{ b }\OperatorTok{*}\NormalTok{ g⁻¹ }\OperatorTok{:=}
\ControlFlowTok{begin}
\NormalTok{    exact (b, by simp),}
\ControlFlowTok{end}

\NormalTok{theorem exercise\_11\_2\_13 (a b }\OperatorTok{:}\NormalTok{ Z) }\OperatorTok{:}
\NormalTok{    (of\_int a }\OperatorTok{:}\NormalTok{ gaussian\_int) }\OperatorTok{|}\NormalTok{ of\_int b }\OperatorTok{→}\NormalTok{ a }\OperatorTok{|}\NormalTok{ b }\OperatorTok{:=}
\ControlFlowTok{begin}
\NormalTok{    contrapose,}
\NormalTok{    simp,}
\ControlFlowTok{end}

\NormalTok{theorem exercise\_1\_1\_17 \{G }\OperatorTok{:} \DataTypeTok{Type}\OperatorTok{*}\NormalTok{\} [group G] \{x }\OperatorTok{:}\NormalTok{ G\} \{n }\OperatorTok{:}\NormalTok{ N\}}
\NormalTok{    (hxn}\OperatorTok{:}\NormalTok{ order\_of x }\OperatorTok{=}\NormalTok{ n) }\OperatorTok{:}
\NormalTok{    x⁻¹ }\OperatorTok{=}\NormalTok{ x }\OperatorTok{\^{}}\NormalTok{ (n }\OperatorTok{{-}} \FloatTok{1} \OperatorTok{:}\NormalTok{ Z) }\OperatorTok{:=}
\ControlFlowTok{begin}
\NormalTok{    rw zpow\_sub\_one,}
\NormalTok{    simp,}
\NormalTok{    rw [}\OperatorTok{←}\NormalTok{ hxn, pow\_order\_of\_eq\_one],}
\ControlFlowTok{end}

\NormalTok{theorem exercise\_3\_1\_22b \{G }\OperatorTok{:} \DataTypeTok{Type}\OperatorTok{*}\NormalTok{\} [group G] (I }\OperatorTok{:} \DataTypeTok{Type}\OperatorTok{*}\NormalTok{)}
\NormalTok{    (H }\OperatorTok{:}\NormalTok{ I }\OperatorTok{→}\NormalTok{ subgroup G) (hH }\OperatorTok{:}\NormalTok{ ∀ i }\OperatorTok{:}\NormalTok{ I, subgroup.normal (H i)) }\OperatorTok{:}
\NormalTok{    subgroup.normal (n (i }\OperatorTok{:}\NormalTok{ I), H i) }\OperatorTok{:=}
\ControlFlowTok{begin}
\NormalTok{    rw infi,}
\NormalTok{    rw }\OperatorTok{←}\NormalTok{set.image\_univ,}
\NormalTok{    rw Inf\_image,}
\NormalTok{    simp [hH],}
\NormalTok{    haveI }\OperatorTok{:=}\NormalTok{ λ i, (H i).normal,}
\NormalTok{    split,}
\NormalTok{    intros x hx g,}
\NormalTok{    rw subgroup.mem\_infi at hx }\OperatorTok{↦}\NormalTok{,}
\NormalTok{    intro i,}
\NormalTok{    apply (hH i).conj\_mem \_ (hx i),}
\ControlFlowTok{end}

\NormalTok{theorem exercise\_3\_4\_5a \{G }\OperatorTok{:} \DataTypeTok{Type}\OperatorTok{*}\NormalTok{\} [group G]}
\NormalTok{    (H }\OperatorTok{:}\NormalTok{ subgroup G) [is\_solvable G] }\OperatorTok{:}\NormalTok{ is\_solvable H }\OperatorTok{:=}
\ControlFlowTok{begin}
\NormalTok{    apply\_instance,}
\ControlFlowTok{end} 
\end{Highlighting}
\end{Shaded}

Figure C: Examples of new proofs discovered by ReProver on ProofNet
{[}29{]}.

data, including file dependencies, ASTs, proof states, tactics, and
premise information. In addition, it enables the model to interact with
Lean 4 through tactics, in the same way as Lean 3 (Sec. 4).

Similar to constructing the Lean 3 version of LeanDojo Benchmark, we
extract data from the commit 3ce43c18f614b76e161f911b75a3e1ef641620ff of
mathlib4 released on October 21, 2023. The resulting dataset is named
LeanDojo Benchmark 4. It is released under the CC BY 2.0 license and
hosted on zenodo.org with DOI ``10.5281/zenodo.8040109''. LeanDojo
Benchmark 4 consists of 102,514 theorems/proofs, 213,067 tactics, and
152,695 premises. We use 2,000 theorems for

\begin{Shaded}
\begin{Highlighting}[]
\NormalTok{theorem exercise\_13\_6\_10 \{K : Type*\} [field K] [fintype K*] :}
\NormalTok{    ∏ (x : K*), x = {-}1 :=}
\NormalTok{begin}
\NormalTok{    exact finite\_field.prod\_univ\_units\_id\_eq\_neg\_one,}
\NormalTok{end}

\NormalTok{theorem exercise\_1\_17}
\NormalTok{    (n : N)}
\NormalTok{    (x y : euclidean\_space R (fin n)) {-}{-} R\^{}n}
\NormalTok{    : ||x + y|\^{}2 + ||x {-} y|\^{}2 = 2*||x|\^{}2 + 2*||y|\^{}2 :=}
\NormalTok{begin}
\NormalTok{    rw [norm\_add\_sq\_real, norm\_sub\_pow\_two\_real],}
\NormalTok{    ring,}
\NormalTok{end}

\NormalTok{theorem exercise\_2\_25 \{K : Type*\} [metric\_space K] [compact\_space K] :}
\NormalTok{    ∃ (B : set (set K)), set.countable B ∧ is\_topological\_basis B :=}
\NormalTok{begin}
\NormalTok{    rcases exists\_countable\_basis K with (B, hBc, hB),}
\NormalTok{    exact (B, hBc, hB.2),}
\NormalTok{end }
\end{Highlighting}
\end{Shaded}

Figure D: Three new proofs discovered by ReProver on ProofNet {[}29{]}
that cannot be found by a baseline without premise retrieval. All of the
three proofs rely on premises:
``finite\_field.prod\_univ\_units\_id\_eq\_neg\_one''\\
,``norm\_add\_sq\_real'',``norm\_sub\_pow\_two\_real'', and
``exists\_countable\_basis''.

validation, 2,000 theorems for testing, and the rest for training.
LeanDojo Benchmark 4 also has two different data splits: random and
novel\_premises.

We use LeanDojo Benchmark 4 to train and evaluate our method. The model
architectures and experimental details are the same as those in Sec. 6.
Results on premise selection are in Table B, and results on theorem
proving are in Table C.

Table B: Premise selection testing performance on LeanDojo Benchmark 4
(Lean 3 results in Table 1). We train and evaluate two models
independently using different data splits (random and novel\_premises).
R@k is the recall for the top k retrieved premises, and MRR is the mean
reciprocal rank metric.

Method

random

novel\_premises

R@1

R@10

MRR

R@1

R@10

MRR

Ours

12.8

34.7

0.29

9.8

32.1

0.24

Table C: Theorem proving Pass@1 (\%) on the testing data of LeanDojo
Benchmark 4 (Lean 3 results in Table 2).

Method

random

novel\_premises

ReProver

48.6

19.9

W/o retrieval

44.5

16.2

\section{E ChatGPT Plugin for Theorem
Proving}\label{e-chatgpt-plugin-for-theorem-proving}

LeanDojo provides a general tool for interacting with Lean
programmatically. As a demo of how it might bridge LLMs and theorem
proving, we build a ChatGPT plugin {[}101{]} enabling ChatGPT to prove
theorems by interacting with Lean through LeanDojo. Plugin developers
can wrap any software

as a web service and describe its APIs to ChatGPT. Then, ChatGPT can
automatically call the APIs and incorporate the results into the
response to the user. Below is a summary of our API description
corresponding to the interface in Sec. 4.

Title: Lean

\begin{Shaded}
\begin{Highlighting}[]
\NormalTok{Description: Plugin for proving user{-}specified theorems automatically by interacting with Lean. The user enters information of how to find a theorem (e.g., theorem name and file path). Based on the user\textquotesingle{}s input, ChatGPT first initializes the proof search with the given theorem as the initial state. Then, ChatGPT will first explain the choice for the next tactic step using LaTeX and run that tactic step to the state. If the current state is not promising, ChatGPT can backtrack to previous states by decrementing the "state\_id" parameter. If applying tactics to the current state specified by the "state\_id" parameter returns an error message, ChatGPT should explain the error, and if repetitive errors occur, ChatGPT should decrement the "state\_id" parameter and try a different approach on a previous state. The theorem is successfully proved if there are no unsolved goals in the current state. }
\end{Highlighting}
\end{Shaded}

Endpoints :

\begin{Shaded}
\begin{Highlighting}[]
\NormalTok{initialize\_proof\_search: Given the theorem name }\KeywordTok{and} \BuiltInTok{file}\NormalTok{ path of a Lean theorem, initialize the proof search. The response includes the initial state }\KeywordTok{and}\NormalTok{ its state ID.}
\NormalTok{Args:}
\NormalTok{    theorem\_name (string): The name of the target theorem to prove.}
\NormalTok{    theorem\_file\_path (string): The }\BuiltInTok{file}\NormalTok{ path of the target theorem.}
\NormalTok{run\_tactic: Run a tactic on a state (specified by its state ID), assuming the proof search has been initialized }\KeywordTok{and}\NormalTok{ some state }\KeywordTok{is}\NormalTok{ available. The response }\KeywordTok{is}\NormalTok{ either the }\BuiltInTok{next}\NormalTok{ state }\KeywordTok{and}\NormalTok{ its state ID }\KeywordTok{or}\NormalTok{ an error message, }\KeywordTok{in}\NormalTok{ which ChatGPT should explain the error }\KeywordTok{and}\NormalTok{ consider decrementing the }\StringTok{"state\_id"}\NormalTok{.}
\NormalTok{Args:}
\NormalTok{    state\_id (string): The ID of the state on which to run the tactic.}
\NormalTok{    tactic (string): The tactic to run on a state (specified by its state ID), assuming the proof search has been initialized. }
\end{Highlighting}
\end{Shaded}

After exposing the APIs to ChatGPT, we can ask it to prove theorems by
specifying the theorem's name and path in any public Lean repo on
GitHub. Fig. E--L show an example with the GPT-3.5 version of ChatGPT.
And Fig. M--O are the same example with the GPT-4 version. The captions
provide detailed step-by-step explanations.

We highlight a few key strengths of ChatGPT observed in multiple
examples we evaluated. First, unlike specialized methods for theorem
proving (this paper and its prior works), ChatGPT interleaved informal
mathematics with formal proof steps. This resembles how humans interact
with proof assistants and opens up new avenues for integrating natural
language and formal theorem proving. Second, ChatGPT demonstrated
impressive capability in explaining error messages from Lean that are
quite opaque even to humans. It was able to incorporate the error
message to refine its proof strategy. Last, ChatGPT's behavior is more
steerable than specialized provers. In Fig. E, we simply gave it the
theorem to prove, but we could also provide more detailed instructions.
For example, we

could say: ``Please describe a high-level proof plan before trying any
tactic.'' This kind of steerability enables future research on prompt
engineering for theorem proving, and we have already seen initial
benefits in an ongoing work named Sagredo. \(^{13}\)

However, these strengths by no means imply ChatGPT can already solve
theorem proving. In fact, it failed to find a proof for most theorems we
tried. Hallucination was common. In Fig. L, ChatGPT falsely asserted the
theorem was proved, while we knew it was not, by looking at LeanDojo's
response. This demonstrates the value of theorem proving as a rigorous
benchmark for addressing LLMs' hallucination problem. Another key
limitation of ChatGPT was its inability to search systematically in a
large space. We frequently found it stuck to an unpromising path when
the correct solution could be found by backtracking and exploring
alternative paths. This behavior is consistent with the general
observation that LLMs are weak at search and planning. Addressing this
weakness is an active area of research {[}102{]}.

We emphasize a few caveats about our study of theorem proving with
ChatGPT. First, data contamination is likely. Many theorems we evaluated
have been publicly available on GitHub before ChatGPT's data cutoff
date. Therefore, ChatGPT may have seen them in training. Second, our
study is exploratory. A more detailed and quantitative study is needed
to characterize ChatGPT's capability in theorem proving. Such a study
with ChatGPT plugins is challenging, as plugins currently only support
interaction through the browser. Also, OpenAI has taken measures to
block automated access by bots. Using humans may be an option, but that
is beyond the scope of this paper.

\textit{[Image omitted: images/abe77537d5072a8e387298a0393186aa6fdbdf719e50ea24ab1d3ca9b7688e2d.jpg]}

text\_image

I want you to prove a theorem in Lean. The theorem's name is
\texttt{hello\_world}, and it is defined in the file
\texttt{src/example.lean} in
\texttt{https://github.com/yangky11/lean-example}.\texttt{Used\ LeanDojo\ LOCALHOST\ REQUEST\ TO\ LEANDOJO\ \{\ \ \ \ \ "theorem\_name":\ "hello\_world",\ \ \ \ \ "theorem\_file\_path":\ "src/example.lean"\ \}\ RESPONSE\ FROM\ LEANDOJO\ \{\ \ \ \ \ "state\_id":\ 0,\ \ \ \ \ "state":\ "a\ b\ c\ :\ N\textbackslash{}n←\ a\ +\ b\ +\ c\ =\ a\ +\ c\ +\ b"\ \}\ The\ theorem}hello\_world\texttt{states\ that\ for\ three\ arbitrary\ complex\ numbers}a\texttt{,}b\texttt{,\ and}c\texttt{,\ the\ sum}a
+ b + c\texttt{is\ equal\ to\ the\ sum}a + c + b`. In other words, it
asserts the commutativity of addition for complex numbers. The goal is
to prove this statement: lean a b c : C ← a + b + c = a + c + b Copy
code

Figure E: (ChatGPT-3.5, 1/8) After receiving the theorem to prove,
ChatGPT first called ``initialize'', which returned the initial state.
Then it tried to interpret the theorem in natural language. Note that it
made a mistake here. The theorem was about natural numbers (N), not
complex numbers (C).

\textit{[Image omitted: images/e26886f293f92df197e94f1007e0e96707b0a53a8c7c2a57b325d7f24912002d.jpg]}\\
Figure F: (ChatGPT-3.5, 2/8) ChatGPT tried to rewrite the goal using the
lemma ``b + c = c + b''. This was a reasonable but incorrect move. After
receiving the error message from Lean, ChatGPT explained the error in
natural language. Here the explanation is quite accurate, which is
impressive given that the error message looks opaque to anyone not
familiar with Lean.

\section{F Limitations and Future
Work}\label{f-limitations-and-future-work}

Our work is one step toward unlocking the potential of LLMs for
generating verifiable formal proofs, and we see abundant space for
future exploration. A learning-based prover is a complex system
consisting of multiple components: data extraction, interaction with
proof assistants, model training, and proof search. While navigating the
design space spanned by various components, we err on the side of
simplicity and efficiency, instead of pushing performance to the limit.
This helps us deliver a reliable, open, and accessible system, laying
the foundation for further research. There are many directions in which
the system can be improved, and we discuss a few of them here. \(^{14}\)

Stronger LLMs. Our backbone model, ByT5 {[}44{]}, was published in 2021
and has 299M parameters, which is not very large by today's standard.
Recently, there have been a plethora of open-source LLMs demonstrating
strong capabilities in writing code, e.g., CodeGen {[}103{]}, StarCoder
{[}94{]}, and CodeGeeX {[}104{]}. We are excited to see how they might
impact theorem proving and, more generally, how far we can go by pushing
the limit of the model/data scale.

ByT5's tokenizer-free nature helps us sidestep the difficulty with
pretrained tokenizers that may not work well for Lean's Unicode-rich
code. However, treating texts as raw bytes makes the sequence length
much longer than necessary. Long sequences harm efficiency, as
Transformers scale quadratically w.r.t. the sequence length, which may
become a bigger problem when we further scale up the model. To solve the
issue, it might be helpful to pretrain a customized tokenizer or adopt
more advanced tokenizer-free models such as MegaByte {[}105{]}.

Our ReProver model is based on the pretraining-finetuning paradigm.
Recent work on instruction-following LLMs such as GPT-4 {[}27{]} has led
to successes in many applications by prompting the model without any
finetuning. Our preliminary results show that GPT-4 and ChatGPT
(Appendix C.2

\textit{[Image omitted: images/a839eb2f7bcb042bd0a5e8580d2bdeabf22d8ffa43152eec3b5d26c58a7c8047.jpg]}\\
Figure G: (ChatGPT-3.5, 3/8) Then it tried to prove the theorem using
``ring''. This was another good move. The ring tactic can prove this
theorem, but Lean could not find it since it was not imported into the
current file. Again, ChatGPT was able to interpret the error message
correctly and concluded that ring was not available. Next, it tried
another tactic but failed again.

and E) cannot solve theorem proving out of the box and are currently far
behind finetuned models. However, the way we prompt these models is
quite naive, and better strategies, such as Tree of Thoughts {[}102{]},
may lead to further improvements. We consider theorem proving as a
promising task for studying LLMs' capabilities in planning and search.

Improving Premise Retrieval. ReProver uses DPR {[}26{]} to retrieve
premises and fuses them with the current proof state by concatenation.
This architecture is simple and effective but does not scale to a large
number of retrieved premises. With a length limit of 2,300 tokens, we
can fit only 10--15 premises into the input of the tactic generator. To
mitigate the problem, we may need an architecture that fuses the
retrieved premises in the hidden space, e.g., Fusion-in-Decoder
{[}106{]}.

\textit{[Image omitted: images/17ceeb4f009117efe28acd04aaa56f6f87d2274bafe234ce3eb62d403b7f7797.jpg]}\\
Figure H: (ChatGPT-3.5, 4/8) ChatGPT made another two failed attempts.
Here, the second attempt had the same problem as in Fig. E (``+'' is
left associative).

In addition, one can also switch from DPR to radically different
retrieval architectures. For example, generative retrieval
{[}107--109{]} is a recent class of models performing retrieval by
directly predicting the document IDs, which could be the premise names
in our task.

Limitations of Imitating Human-Written Proofs. Human-written proofs
extracted by LeanDojo provide valuable data for training the prover.
However, we have also observed limitations of using them as the sole
training target:

First, they are relatively scarce for today's data-hungry LLMs. LeanDojo
Benchmark has 98,734 proofs, covering a large portion of available data
in Lean (as of October 2023). The number of proofs

\textit{[Image omitted: images/eff5557a7e17050a32cc21639918e9ae732aa6883a2309be97069ddd69e2226e.jpg]}\\
Figure I: (ChatGPT-3.5, 5/8) ChatGPT tried to simplify the goal using
the simp tactic but failed. Then it tried the ``rw add\_comm a b''
tactic, which successfully rewrote the goal from ``a + b + c = a + c +
b'' to ``b + a + c = a + c + b''.

in other proof assistants has the same order of magnitude (tens or
hundreds of thousands). Due to limited data, we cannot constantly
improve the performance simply by scaling up the model size. Second,
theorem proving in proof assistants is an interactive process, but the
proof only captures the final successful trajectory. Without the
intermediate history of trial and error, it can be quite opaque how
final proofs are derived. Therefore, tactics in human-written proofs can
be difficult for the model to learn from. Third, models trained on
proofs in one project often struggle to generalize to theorems in new
domains {[}9, 20{]}, e.g., from mathlib to MiniF2F and ProofNet
(Appendix C.4).

\textit{[Image omitted: images/7f79ea9166d2fbe9e8729c157aa72af66620066b8ca4925d6b365b94d814e99d.jpg]}

text\_image

Now, we can use the \texttt{tw} tactic again to apply the commutativity
of addition to the subterm \texttt{a\ +\ c} in the goal. This will allow
us to rewrite the goal as \texttt{b\ +\ a\ +\ c\ =\ c\ +\ a\ +\ b}.
Let's try this approach. Used LeanDojo LOCALHOST REQUEST TO LEANDOJO \{
``state\_id'': 1, ``tactic'': ``rw add\_comm a c'' \} RESPONSE FROM
LEANDOJO \{``state\_id': 2, `state': `a b c : N\n-b + a + c = c + a +
b', `proof\_finished': False\}'' After applying the commutativity of
addition to the subterm \texttt{a\ +\ c}, we have successfully rewritten
the goal as \texttt{b\ +\ a\ +\ c\ =\ c\ +\ a\ +\ b}. The updated goal
is: lean a b c : C ← b + a + c = c + a + b

Figure J: (ChatGPT-3.5, 6/8) Then it successfully performed another
rewrite from ``b + a + c = a + c + b'' to ``b + a + c = c + a + b''.

To overcome these limitations, existing work has explored learning from
auxiliary data or data collected via online interaction with the proof
assistant. For example, Proof Artifact Co-Training (PACT) co-trains the
tactic generator on nine auxiliary tasks, such as predicting types and
theorem names {[}16{]}. MetaGen {[}38{]} trains a neural network to
generate synthetic theorems/proofs as training data in the Metamath
proof assistant {[}110{]}. Polu et al.~{[}19{]} and Lample et
al.~{[}17{]} improve the prover by training it on successful proofs
found by itself. Incorporating these techniques into our system may lead
to substantial improvements.

\textit{[Image omitted: images/79d14566c97b521527fa9aeb38b5ba1ba89a4f180a7f0153644a4b2910283255.jpg]}\\
Figure K: (ChatGPT-3.5, 7/8) Another rewrite tactic, transforming the
goal to ``a + b + c = c + a + b''.

\textit{[Image omitted: images/1e7430e0e7299d5ba92c050e68ea635de8a60721d2922bb4c9640c48848efc7e.jpg]}

text\_image

Finally, we can use the
\texttt{tw\textquotesingle{}\ tactic\ one\ last\ time\ to\ apply\ the\ commutativity\ of\ addition\ to\ the\ subterm}c
+
a\texttt{in\ the\ goal.\ This\ will\ allow\ us\ to\ rewrite\ the\ goal\ as}a
+ b + c = a + c +
b\texttt{,\ which\ matches\ the\ original\ goal\ and\ completes\ the\ proof.\ Let\textquotesingle{}s\ finish\ the\ proof.\ Used\ LeanDojo\ \ \ \ LOCALHOST\ REQUEST\ TO\ LEANDOJO\ \{\ \ \ \ \ "state\_id":\ 3,\ \ \ \ \ "tactic":\ "rw\ add\_comm\ c\ a"\ \}\ RESPONSE\ FROM\ LEANDOJO\ \{"state\_id":\ 4,\ \textquotesingle{}state\textquotesingle{}:\ \textquotesingle{}a\ b\ c\ :\ N\textbackslash{}ni-\ a\ +\ b\ +\ c\ =\ a\ +\ c\ +\ b\textquotesingle{},\ \textquotesingle{}proof\_finished\textquotesingle{}:\ False\}"\ It\ looks\ like\ the\ goal}a
+ b + c = a + c +
b\texttt{remains\ unchanged\ after\ applying\ the\ commutativity\ of\ addition\ to\ the\ subterm}c
+
a\texttt{.\ This\ is\ because\ the\ goal\ is\ already\ in\ the\ desired\ form,\ and\ the\ proof\ is\ complete.\ Since\ the\ goal\ matches\ the\ original\ statement\ of\ the\ theorem,\ we\ have\ successfully\ proved\ the\ theorem}hello\_world\texttt{.\ The\ proof\ demonstrates\ that\ for\ three\ arbitrary\ complex\ numbers}a\texttt{,}b\texttt{,\ and}c\texttt{,\ the\ sum}a
+ b + c\texttt{is\ equal\ to\ the\ sum}a + c + b`, which confirms the
commutativity of addition for complex numbers.

Figure L: (ChatGPT-3.5, 8/8) It performed one last rewrite, which
brought the goal back to the original goal ``a + b + c = a + c + b''. At
this point, ChatGPT went back to the starting point without making
meaningful progress. However, it asserted that the theorem had been
successfully proved, even though the response from LeanDojo clearly
indicated ``proof\_finished: False''.

\textit{[Image omitted: images/c9e0323bca6dc5150df5952bf8b0c22e06d56df79f042947d69f2365bbf80ccc.jpg]}

text\_image

I want you to prove a theorem in Lean. The theorem's name is
\texttt{hello\_world}, and it is defined in the file
\texttt{src/example.lean} in
\texttt{https://github.com/yangky11/lean-example}. Used LeanDojo
LOCALHOST REQUEST TO LEANDOJO \{ ``theorem\_name'': ``hello\_world'',
``theorem\_file\_path'': ``src/example.lean'' \} RESPONSE FROM LEANDOJO
\{ ``state\_id'': 0, ``state'': ``a b c : N\n+ a + b + c = a + c + b''
\} The theorem \texttt{hello\_world} is stated as follows: For any three
integers \texttt{a}, \texttt{b}, and \texttt{c}, the sum
\texttt{a\ +\ b\ +\ c} is equal to the sum \texttt{a\ +\ c\ +\ b} In
other words, it asserts the commutativity of addition for integers,
which is a well-known property of integers. Let's start the proof. The
first step is to apply the \texttt{rw} (rewrite) tactic with the
\texttt{add\_assoc} lemma, which states that
\texttt{(a\ +\ b)\ +\ c\ =\ a\ +\ (b\ +\ c)}. This will allow us to
rearrange the terms in the sum. Used LeanDojo LOCALHOST REQUEST TO
LEANDOJO \{ ``state\_id'': 0, ``tactic'': ``rw add\_assoc'' \} RESPONSE
FROM LEANDOJO \{\{\texttt{state\_id\textquotesingle{}:\ 1,}state':
\texttt{a\ b\ c\ :\ N\textbackslash{}n+\ a\ +\ (b\ +\ c)\ =\ a\ +\ c\ +\ b},
\texttt{proof\_finished\textquotesingle{}:\ False\}"\ \}\ After\ applying\ the}add\_assoc\texttt{lemma,\ we\ have\ the\ new\ goal:\ For\ any\ three\ integers}a\texttt{,}b\texttt{,\ and}c\texttt{,\ the\ sum}a
+ (b + c)\texttt{is\ equal\ to\ the\ sum}a + c + b`.

Figure M: (ChatGPT-4, 1/3) The same theorem using the GPT-4 version of
ChatGPT. Despite a minor error when interpreting the theorem (natural
numbers, not integers), it generated a sensible tactic, rewriting the
goal to ``a + (b + c) = a + c + b''.

\textit{[Image omitted: images/d9666f69e0cafc360f39a02a9507aeda92f7faa92dd4cbb2e2db4970aa05b807.jpg]}\\
Figure N: (ChatGPT-4, 2/3) A good move, rewriting the goal to ``a + (c +
b) = a + c + b''.

\textit{[Image omitted: images/2e054b9b3b58143c44da59001cbd1aaad04bec03336ada330bc018d44435a5cc.jpg]}\\
Figure O: (ChatGPT-4, 3/3) It tried the refl tactic but failed. refl
requires the goal to be an equation whose both sides are identical up to
trivial transformations. However, ``a + (c + b) = a + c + b'' is not
trivial since ``+'' is left associative. ChatGPT was able to interpret
this error accurately and finish the proof using the correct premise
``add\_assoc''.

\end{document}
